\specialchapter{警告}

这次重新印行的流形结果分册不加改动地重印了此前出版的著作（第 XXXIII 和 XXXXI 分册），并将它们汇编为一卷。

Nancago，1982 年秋

N. BOURBAKI

\part{第 1 至 7 段（97 页）}

\specialchapter{引言}

本分册汇集了可微流形（在实数域上）理论和解析流形（在完备非离散赋值域上）理论的基本概念与主要结果。不含证明。

假定已知前六卷的定义和结果。

\specialchapter{记号与约定（第 1 至 7 段）}

\specialsection{基域}

在本分册中，字母 K 表示以下三者之一：实数的赋值域 R、复数的赋值域 C，或带超度量绝对值的完备交换非离散赋值域（Top. Gen.，第 IX 章，第 3 版，第 3 节）。

\specialsection{拓扑向量空间}

当 K 不同于 R 或 C 时，我们把 K 上向量空间 F 的半范数（分别地，范数）称为 F 上的超半范数（分别地，作为范数的超半范数）（Esp. Vect. Top.，第 II 章，第 2 版，第 1 节）。若 γ 是 F 上的半范数且 x ∈ F，有时也写作 \textbar\textbar x\textbar\textbar γ 而不写 γ(x)。

我们把拓扑可由单个范数定义的 K 上拓扑向量空间称为可赋范空间，把完备的可赋范空间 \(^{1}\) 称为巴拿赫空间。我们把拓扑可由一族半范数定义的 K 上拓扑向量空间称为多范数空间；当 K = R 或 C 时，这一概念与局部凸空间的概念一致。

若 E 和 F 是多范数空间，我们用 \(\mathcal{L}_{m}(\mathrm{E};\mathrm{F})\) 表示从 \(m\) 到 F 的连续 \(\mathbf{E}^{m}\) 线性映射空间，并赋予其在 \(\mathbf{E}^{m}\) 的有界子集上一致收敛的拓扑；它是一个多范数空间。此外，若 E 可赋范，\(\gamma\) 是 F 上的连续半范数，且 \(u\in\mathcal{L}_{m}(\mathrm{E};\mathrm{F})\)，则以 \(\|u\|_{\gamma}\) 表示满足下列条件的实数 \(a\geqslant0\) 的下确界：

\[
\left\| u (x _ {1}, \dots , x _ {m}) \right\| _ {\gamma} \leqslant a \| x _ {1} \| \dots \| x _ {m} \|
\]

对 E 中所有的 \(x_{1}, \ldots, x_{m}\) 均成立。

\specialsection{多重指标}

设 I 为一集合。对于属于 \(\alpha = (\alpha_{i})\) 的 \(\beta = (\beta_{i})\) 和 \(\mathbf{N}^{(I)}\)，置（参见 Alg.，第 III 章，第 \(3^{\mathrm{e}}\) 版，章首）：

\[
| \alpha | = \sum_ {i \in I} \alpha_ {i}
\]

\[
\alpha ! = \prod_ {i \in I} \alpha_ {i}!
\]

\[
\alpha + \beta = (\alpha_ {i} + \beta_ {i}) _ {i \in I}
\]

\[
((\alpha , \beta)) = \prod_ {i \in I} \frac {(\alpha_ {i} + \beta_ {i}) !}{\alpha_ {i} ! \beta_ {i} !}.
\]

当对所有 \(\alpha \leqslant \beta\) 都有 \(\alpha_{i} \leqslant \beta_{i}\) 时，置 \(i \in I\)；这等价于存在 \(\gamma \in \mathbf{N}^{(I)}\) 使得 \(\beta = \alpha + \gamma\)；此时该元素 \(\gamma\) 唯一，记作 \(\beta - \alpha\)。

若 \(\mathbf{x} = (x_i)_{i \in I}\) 是环 A 中两两可交换的元素族，且 A 具有单位元 1，则置：

\[
\mathbf {x} ^ {\alpha} = \prod_ {i \in \mathbf {I}} x _ {i} ^ {\alpha_ {i}} \quad (\text { with   the   convention   that } x ^ {0} = 1).
\]

有 \(\mathbf{x}^{\alpha +\beta} = \mathbf{x}^{\alpha}\cdot \mathbf{x}^{\beta}\)。

对于 \(i\in I\)，用 \(\varepsilon_{i}\) 表示 \(\mathbf{N}^{(I)}\) 中除指标为 \(i\) 的坐标等于 1 外其余坐标全为零的元素。对于 \(\alpha=\sum_{i\in I}\alpha_{i}\varepsilon_{i}\) 中每个元素 \(\alpha=(\alpha_{i})\)，都有 \(\mathbf{N}^{(I)}\)。

\specialsection{映射的限制}

若 \(f\) 是从集合 \(A\) 到集合 \(B\) 的映射，且 \(C\) 是 \(A\) 的子集，则用 \(f|C\) 表示 \(f\) 对 \(C\) 的限制。

\specialsection{元素 ∞ 与 ω}

在整数集上添入两个记作 ∞ 和 ω 的元素。整数之间的序关系按如下约定扩充：

\[
n <   \infty <   \omega \quad \text { for   every   integer } n.
\]

置：

\[
\infty + n = n + \infty = \infty \quad \text { and } \quad \omega + n = n + \omega = \omega , \quad n \in \mathbf {Z}.
\]

字母 \(r\) 表示一个 \(\geqslant 1\) 的整数，或元素 \(\infty\)、\(\omega\) 之一。当 \(K \neq R\) 时，字母 \(r\) 总表示元素 \(\omega\)。当 \(r = \infty\)（分别地，当 \(r = \omega\)）时，对每个整数 \(r - k = r\) 置 \(k \geqslant 1\)。

\section{可微函数}

本节中，字母 E 表示 K 上的可赋范拓扑向量空间；字母 F 表示 K 上的分离多范数空间。

\subsection{点处两个函数的接触阶}

1.1.1. 设 X 为拓扑空间，\(\theta\) 为 X 中一点 \(x_0\) 的某邻域上定义的正数值函数。若函数 \(f\) 在 \(x_0\) 的某邻域上定义且取值于 F，并满足下列条件，则称 f 在 \(\theta\) 处相对于 \(x_0\) 可忽略：

对每个 \(\varepsilon > 0\) 以及 F 上每个连续半范数 \(\gamma\)，存在 \(F\) 的一个邻域 \(V\)\(x_0\)，使 \(f\) 和 \(\theta\) 在其上有定义，并且

\[
\| f (x) \| _ {\gamma} \leqslant \varepsilon \theta (x) \quad \text { for   every } x \in \mathbf {V}.
\]

要使 \(f\) 相对于 \(\theta\) 可忽略，只需对定义 F 的拓扑的一族半范数 \(\gamma\) 满足此条件即可。\(f\) 在 \(\theta\) 处相对于 \(x_{0}\) 是否可忽略，只取决于 \(f\) 和 \(\theta\) 在 \(x_{0}\) 处的芽。用 \(o_{x_{0}}(\theta)\)（当 \(o(\theta)\) 不会引起歧义时也写作 \(x_{0}\)）表示在 \(x_{0}\) 处相对于 \(\theta\) 可忽略的函数在 \(x_{0}\) 处的芽的集合：它是取值于 F 的映射在 \(x_{0}\) 处的芽空间的向量子空间。若 \(f\) 相对于 \(\theta\) 可忽略，按记号滥用，我们写 \(f \in o_{x_{0}}(\theta)\)，或者当 \(f(x) \in o(\theta(x))\) 趋于 \(x\) 时写 \(x_{0}\)。

若 \(f\) 和 \(g\) 是从 \(x_{0}\) 的某邻域到 \(F\) 的两个映射，且 \(f \equiv g \mod o(\theta)\) 相对于 \(f - g\) 可忽略，则也写作 f  g  o(\(\theta\))。

假设 K 等于 R 或 C，且 \(x_{0}\) 属于 X 中 \(\theta\) 有定义且非零的点集 Y 的闭包。此时关系 \(f \in o(\theta)\) 意味着：x 在 Y 中趋于 \(f(x)/\theta(x)\) 时，\(x_{0}\) 趋于 0，并且 \(\theta(x) = 0\) 蕴含 \(f(x) = 0\)。

1.1.2. 设 \(f\) 和 \(g\) 是在 \(F\) 中一点 \(x_0\) 的某邻域上定义且取值于 \(E\) 的两个函数。若 \(m\) 是正整数，且有下式，则称 \(f\) 和 \(g\) 在 \(\geqslant m\) 处的接触阶 \(x_0\)：

\[
f (x) - g (x) \in o (\| x - x _ {0} \| ^ {m}) \quad \text { for } x \text { tending   toward } x _ {0}
\]

无论选取何种范数来定义 E 的拓扑，均应如此。为此，只需对定义 E 的拓扑的某个范数验证上述关系即可。若如此，则有 \(f(x_0) = g(x_0)\)。

若 \(f\) 和 \(g\) 在 \(x_{0}\) 处取相同值，则称 \(f\) 和 \(g\) 在 \(x_{0}\) 处的接触阶为满足 \(+\infty\) 和 \(m\) 在 \(f\) 处接触阶 \(g\) 的整数 \(\geqslant m\) 的上确界（有限或等于 \(x_{0}\)）。

1.1.3. \(f\) 和 \(g\) 在 \(x_{0}\) 处的接触阶只取决于它们在点 \(x_{0}\) 处的芽。因此，可以谈论从 E 到 F 的两个映射芽 \(\varphi\) 和 \(\psi\) 在点 \(x_{0}\) 处的接触阶。关系 \(\langle\varphi\) 与 \(\psi\) 的接触阶 \(\geqslant m\rangle\) 是一个与向量结构相容的等价关系。

\subsection{点处的可微函数}

1.2.1. 设 \(f\) 在 E 中点 \(x_{0}\) 的某邻域上定义且取值于 F。若存在从 E 到 F 的连续仿射函数 \(f\)，使得 \(x_{0}\) 与 \(v\) 在 \(x_{0}\) 处的接触阶 \(\geqslant 1\)，则称 \(f\) 在 x＿0 处可微。此映射 \(v\) 唯一；存在唯一的从 E 到 F 的连续线性映射，记作 \(Df(x_{0})\)，使得：

\[
v (x) = v \left(x _ {0}\right) + D f \left(x _ {0}\right). \left(x - x _ {0}\right).
\]

若在 E 上选取一个范数，则这等价于：

\(f(x_0 + h) \equiv f(x_0) + Df(x_0).h \mod o(\|h\|)\)（其中 \(h\) 趋于 0），

这也可写成如下形式：

\[
\lim _ {h \rightarrow 0, h \neq 0} \frac {\| f (x _ {0} + h) - f (x _ {0}) - D f (x _ {0}) . h \| _ {\gamma}}{\| h \|} = 0
\]

对 F 上每个连续半范数 γ。

称 \(Df(x_{0})\) 中的元素 \(\mathcal{L}(\mathrm{E},\mathrm{F})\) 为 f 在 \(x_{0}\) 处的导数。有时用 \(D_{h}f(x_{0})\) 表示 \(Df(x_{0}).h\)。它是 F 中由下式定义的元素：

\[
\mathrm{D} _ {h} f (x _ {0}) = \lim _ {t \rightarrow 0, t \neq 0} \frac {f (x _ {0} + t h) - f (x _ {0})}{t}.
\]

1.2.2. 若函数 \(f\) 在 \(x_{0}\) 处可微，且对每个定义 E 的拓扑的范数都有下列关系，则称 f 在 \(x_{0}\) 处严格可微：

\[
f (y) - f (z) \equiv D f (x _ {0}). (y - z) \quad \text {   mod   } o (\| y - z \|)
\]

当 \((y, z)\) 在 \((x_0, x_0)\) 中趋于 \(E \times E\) 时成立。为此，只需对定义 \(E\) 的拓扑的某个范数满足此条件即可。此外，若 \(E\) 和 \(F\) 都是赋范空间，则对每个数 \(c > \|Df(x_0)\|\)，存在 \(x_{0}\) 的一个邻域 V，使得对 V 中的 y、z 有 \(\|f(y)-f(z)\|\leqslant c.\|y-z\|\)；这蕴含 f 在 V 中一致连续。

1.2.3. 函数 \(f\) 在 \(x_{0}\) 处可微或严格可微这一性质只取决于 \(f\) 在 \(x_{0}\) 处的芽。在 \(x_{0}\) 处可微的函数芽构成所有芽空间的向量子空间 \(\mathcal{V}\)，且从 \(f\mapsto\mathrm{D}f(x_{0})\) 到 \(\mathcal{V}\) 的映射 \(\mathcal{L}(\mathrm{E};\mathrm{F})\) 是线性的。在 \(x_{0}\) 处严格可微的函数芽构成 \(\mathcal{V}\) 的向量子空间。

1.2.4. 在 \(x_0\) 处可微的函数在 \(x_0\) 处连续。

1.2.5. 当 E = K 时，映射 \(u \mapsto u(1)\) 是从 \(\mathcal{L}(\mathbf{E}; \mathbf{F})\) 到 F 的同构；若函数 f 在 \(x_{0}\) 处可微，则元素

\[
f ^ {\prime} (x _ {0}) = \mathrm{D} f (x _ {0}). 1
\]

正是 Fonct. Var. Réelle，第 I 章，第 1 节，第 6 号注 2 中所给意义下的 \(f\) 在 \(x_0\) 处的导数。

\subsection{可微函数的复合}

1.3.1. 假设 F 可赋范。设 \(x_0 \in E\)、\(y_0 \in F\)，U 是 \(x_0\) 的一个邻域，V 是 \(y_0\) 的一个邻域；最后，设 \(f\) 是从 U 到 V 的映射，在 \(x_0\) 处可微，且 \(f(x_0) = y_0\)。若 \(g\) 是从 V 到分离多范数向量空间 G 的映射，在 \(y_0\) 处可微，则从 U 到 G 的映射 \(g \circ f\) 在 \(x_0\) 处可微，且有：

\[
\mathbf {D} (g \circ f) (x _ {0}) = \mathbf {D} g (y _ {0}) \circ \mathbf {D} f (x _ {0}).\tag{1}
\]

若 \(f\) 和 \(g\) 严格可微，则 \(g \circ f\) 也有同样结论。

1.3.2. 设 \(f\) 是在 E 中一点 \(x_0\)\(E\) 的某邻域上定义且取值于 \(F\) 的映射，并在 \(x_0\) 处可微；若 \(u\) 是从 \(F\) 到分离多范数空间 \(G\) 的连续线性映射，则函数 \(u \circ f\) 在 \(x_0\) 处可微，且有：

\[
\mathrm{D} (u \circ f) (x _ {0}) = u \circ \mathrm{D} f (x _ {0}).\tag{2}
\]

1.3.3. 假设 F 是分离多范数向量空间族 \((F_{i})_{i\in I}\) 的乘积；对 I 中每个 i，设 \(f_{i}\) 是在 E 中一点 \(x_{0}\) 的邻域 U 上定义且取值于 F 的映射，并令 \(f = (f_{i})_{i\in I}\)。f 在 \(x_{0}\) 处可微（分别地，严格可微）的充要条件是所有 \(f_{i}\) 都具有该性质；有：

\[
\mathrm{D} _ {h} f (x _ {0}) = (\mathrm{D} _ {h} f _ {i} (x _ {0})) _ {i \in \mathbf {I}} \quad \text { for   every } h \text { in } E.\tag{3}
\]

1.3.4. 若 E = K，则在公式 (1) 和 (3) 中可用 \(Df(x_{0})\) 替代 \(f'(x_{0})\)，用 \(D_{h}f_{i}(x_{0})\) 替代 \(f_{i}'(x_{0})\)。

\subsection{可微函数的乘积}

1.4.1. 设 \(F_{1}, \ldots, F_{m}\) 为分离多范数空间，u 为从 \(F_{1} \times \cdots \times F_{m}\) 到 F 的连续 m 线性映射。设 U 是 E 中一点 \(x_{0}\) 的邻域，\(f_{i}\) 是从 U 到 \(F_{i}\) 的映射（\(1 \leqslant i \leqslant m\)）。若 \(f_{i}\) 在 \(x_{0}\) 处可微（分别地，严格可微），则 \(u(f_{1}, \ldots, f_{m}) = g\) 也具有该性质，且有：

\[
\mathrm{D} _ {h} g (x _ {0}) = \sum_ {j = 1} ^ {m} u (f _ {1} (x _ {0}), \dots , \mathrm{D} _ {h} f _ {j} (x _ {0}), \dots , f _ {m} (x _ {0})) \quad \text { for } h \text { in } E, \tag {4}
\]

这可简写为：

\[
\mathrm{D} g = \sum_ {j = 1} ^ {m} u (f _ {1}, \dots , \mathrm{D} f _ {j}, \dots , f _ {m}).\tag{5}
\]

特别地，当 \(m = 2\) 时，有：

\[
\mathrm{D} u (f _ {1}, f _ {2}) = u (\mathrm{D} f _ {1}, f _ {2}) + u (f _ {1}, \mathrm{D} f _ {2}).\tag{6}
\]

当 \(m = 1\) 时，得到 1.3.2。

1.4.2. 当 E = K 时，在公式 (4) 至 (6) 中可用 \(g'\) 替代 Dg，用 \(Df_{j}\) 替代 \(f_{j}'\)。

\subsection{隐函数定理的第一类变体}

设 E 和 F 是巴拿赫空间，\(x_{0}\) 是 E 中一点，U 是 \(x_{0}\) 的一个邻域，f 是从 U 到 F 的映射。再设 f 在 \(x_{0}\) 处严格可微。

若 \(Df(x_0)\) 是从 E 到 F 的同构，则存在包含于 U 的 \(U_0\) 的开邻域 \(x_0\)，以及 \(V_0\) 的开邻域 \(f(x_0)\)，使得 \(f|U_0\) 是从 \(U_0\) 到 \(V_0\) 的同胚。逆于 \(g: V_0 \to U_0\) 的映射 \(f|U_0\) 在点 \(f(x_0)\) 处严格可微，且有：

\[
\mathrm{Dg} (f (x _ {0})) = \mathrm{Df} (x _ {0}) ^ {- 1}.
\]

1.5.2. 若 \(Df(x_{0})\) 是从 E 到 F 的满射，则存在包含于 U 的 \(U_{0}\) 的开邻域 \(x_{0}\)，使得 \(f|U_{0}\) 是开映射。

1.5.3. 若 \(Df(x_{0})\) 是单射且像为闭集，则存在包含于 U 的 \(U_{0}\) 的闭邻域 \(x_{0}\)，使得 \(f|U_{0}\) 是从 \(U_{0}\) 到 F 的闭子集的同胚。

\subsection{偏导数}

1.6.1. 设 \(f\) 在 \(U\) 中点 \(x_0\) 的邻域 \(E\) 上定义且取值于 \(F\)。设 \(X\) 是 \(E\) 的向量子空间，\(V\) 是 \(x\) 中满足 \(X\) 的点 \(x_0 + x \in U\) 的集合；对 \(g(x) = f(x_0 + x)\) 置 \(x \in V\)。若 \(f\) 在 0 处有导数，则称 \(X\) 在 \(x_0\) 处有关于 \(g\) 的偏导数；此导数记作 \(D_X f(x_0)\)；它是从 \(X\) 到 \(F\) 的连续线性映射。若 \(f\) 在 \(x_0\) 处可微，则 f 在 \(X\) 处有关于 \(x_0\) 的偏导数，且此偏导数是 \(Df(x_0)\) 对 \(X\) 的限制。

1.6.2. 假设 \(E\) 是有限个赋范向量空间 \(E_{i}\)（\(1 \leqslant i \leqslant n\)）的乘积，并将它们按规范方式识别为 \(E\) 的子空间；设 \(x_{0} = (x_{0}^{1}, \ldots, x_{0}^{n})\) 属于 \(E\)，\(U\) 是 \(x_{0}\) 中 \(E\) 的一个邻域；最后设 \(f\) 是从 \(U\) 到 \(F\) 的映射。若存在，则用 \(D_{i}f(x_{0})\) 表示在点 \(x_{0}^{i}\) 处的映射 \(z_{i} \mapsto f(x_{0}^{1}, \ldots, z_{i}, \ldots, x_{0}^{n})\) 的导数，该映射在 \(x_{0}^{i}\) 中 \(E_{i}\) 的某邻域上定义且取值于 \(F\)。这是 \(\mathcal{L}(E_{i}; F)\) 中的元素，称为 \(f\) 在 \(x_{0}\) 处的第 i 个偏导数。若 \(f\) 在 \(x_{0}\) 处可微，则 \(n\) 个偏导数均存在，并由下式确定 \(Df(x_{0})\)：

\[
\mathrm{D} f (x _ {0}). h = \sum_ {i = 1} ^ {n} \mathrm{D} _ {i} f (x _ {0}). h _ {i} \quad \text { for } h = (h _ {1}, \dots , h _ {n}) \text { in } E.
\]

1.6.3. 更具体地，令 \(E = K^{n}\)。若 f 在 \(x_{0}\) 处的偏导数存在，则用 \(\partial_{i}f(x_{0})\) 表示 F 中的元素 \(D_{i}f(x_{0}).1\)。常使用以下记号；假设已为 \(K^{n}\) 上的坐标函数选定记号，例如 \(u_{i}\) 表示从 \(K^{n}\) 到 K 的第 i 个投影。于是写作：

\[
\frac {\partial f}{\partial u _ {i}} \left(x _ {0}\right) \quad \text { or } \quad \left. \frac {\partial f}{\partial u _ {i}} \right| _ {x = x _ {0}}
\]

而不写 \(\partial_{i}f(x_{0})\)。

1.6.4. 令 \(E = K^{n}\)、\(F = K^{m}\)；设取值于 F 的函数 \(f = (f_{1}, \ldots, f_{m})\) 在 E 中点 \(x_{0}\) 处可微。偏导数 \(a_{ji} = \partial_{i} f_{j}(x_{0})\) 于是存在（它们是 K 中的元素）。由 \(a_{ji}\) 构成的 m 行 n 列矩阵（元素位于指标为 j 的行和指标为 i 的列）称为 f 在 \(x_{0}\) 处的雅可比矩阵；相对于这些空间的规范基，它是从 \(Df(x_{0})\) 到 \(K^{n}\) 的线性映射 \(K^{m}\) 的矩阵。

\subsection{迭代导数}

1.7.1. 设 \(f\) 在 E 中一点 \(x_{0}\) 的某邻域上定义且取值于 F。若 \(f\) 在 \(x_{0}\) 的某邻域上可微，则其导数 Df 是从 \(x_{0}\) 的某邻域到多范数空间 \(\mathcal{L}(\mathrm{E};\mathrm{F})\) 的映射，其中后者是从 E 到 F 的连续线性映射空间。设 \(p\) 是满足 \(\geqslant 2\) 的整数：

若 \(f\)\(p\) 在 \(x_0\) 的某邻域上可微，且其导数 \(f\) 在 \(x_0\) 处 \(Df\) 次可微，则称 \((p - 1)\) 在 \(x_0\) 处 \(p\) 次可微。于是定义 \(f\) 在 \(x_0\) 处的 \(p\) 阶导数；这是由下式定义的从 \(D^p f(x_0)\) 到 F 的连续 \(E^p\) 线性映射 \(F\)：

\[
\mathrm{D} ^ {p} f (x _ {0}). (h _ {1}, \dots , h _ {p}) = (\mathrm{D} (\mathrm{D} ^ {p - 1} f) (x _ {0}). h _ {1}). (h _ {2}, \dots , h _ {p}).
\]

此外置 \(D^{0}f = f\) 且 \(D^{1}f = Df\)。若 \(f\) 在 \(p\) 处 \(x_{0}\) 次可微，且 \(q\)、\(s\) 是满足 \(q + s = p\) 且 \(s > 0\) 的两个整数，则 \(f\) 在 \(q\) 的某邻域上 \(x_{0}\) 次可微，取值于 \(D^{q}f\) 的函数 \(\mathcal{L}_{q}(E; F)\) 在 \(s\) 处 \(x_{0}\) 次可微，并且有：

\[
\mathrm{D} ^ {q + s} f (x _ {0}). (h _ {1}, \dots , h _ {q + s}) = (\mathrm{D} ^ {s} (\mathrm{D} ^ {q} f) (x _ {0}). (h _ {1}, \dots , h _ {s})) \cdot (h _ {s + 1}, \dots , h _ {q + s})
\]

这一关系按记号滥用写成：

\[
\mathbf {D} ^ {q + s} f = \mathbf {D} ^ {s} \mathbf {D} ^ {q} f.
\]

1.7.2. 设 \((\mathbf{E}_{i})_{1\leqslant i\leqslant n}\) 是 E 的闭向量子空间，且 E 是各 \(E_{i}\) 的拓扑直和。于是，在存在时定义从 \(D_{i_{1}}\ldots D_{i_{m}}f\) 的某邻域到 F 的映射 f 的迭代偏导数 \(x_{0}\in E\)；它是从

\[
\mathbf {E} _ {i _ {1}} \times \dots \times \mathbf {E} _ {i _ {m}}
\]

到 F 的连续多线性映射，按整数 \(m \geqslant 1\) 归纳定义如下：若 \(D_{i_{2}} \ldots D_{i_{m}} f(x)\) 在 \(x_{0}\) 的某邻域上存在，且关于 \(E_{i_{1}}\) 有偏导数，则 \(D_{i_{1}} \ldots D_{i_{m}} f(x_{0})\) 由下式给出：

\[
\mathrm{D} _ {i _ {1}} \dots \mathrm{D} _ {i _ {m}} f (x _ {0}). (h _ {1}, \dots , h _ {m}) = (\mathrm{D} _ {i _ {1}} (\mathrm{D} _ {i _ {2}} \dots \mathrm{D} _ {i _ {m}} f) (x _ {0}). h _ {1}). (h _ {2}, \dots , h _ {m})
\]

其中 \(h_{k} \in E_{i_{k}}\)。

若 \(f\) 在 \(m\) 处 \(x_{0}\) 次可微，则偏导数 \(\mathbf{D}_{i_{1}}\ldots\mathbf{D}_{i_{m}}f(x_{0})\) 存在，并等于 \(\mathbf{D}^{m}f(x_{0})\) 对 \(\mathbf{E}_{i_{1}}\times\cdots\times\mathbf{E}_{i_{m}}\) 的子空间 \(\mathbf{E}^{m}\) 的限制。因此，\(\mathbf{D}^{m}f(x_{0})\) 完全由 \(m\) 处 \(x_{0}\) 阶迭代偏导数确定。

1.7.3. 假设 E 是有限维的，且 \((e_{1},\ldots,e_{n})\) 是 E 的一组基。置 \(E_{i}=Ke_{i}\)，并设 f 是从 \(x_{0}\) 的某邻域到 F 的映射。若偏导数 \(D_{i_{1}}\ldots D_{i_{m}}f(x_{0})\)（使用 1.7.2 的记号）存在，则置：

\[
\partial_ {i _ {1}} \dots \partial_ {i _ {m}} f (x _ {0}) = \mathbf {D} _ {i _ {1}} \dots \mathbf {D} _ {i _ {m}} f (x _ {0}). (e _ {i _ {1}}, \dots , e _ {i _ {m}}).
\]

\section{实可微函数}

本节假定 K = R。字母 E 表示 R 上的赋范向量空间；字母 F 表示 R 上的分离局部凸拓扑向量空间。

\subsection{点处的可微函数}

2.1.1. 设 \(f\) 在 E 中一点 \(x_{0}\) 的某邻域上定义且取值于 F。设 \(u\) 是从 E 到 F 的连续线性映射空间 \(\mathcal{L}(\mathrm{E};\mathrm{F})\) 中的元素。\(f\) 在 \(x_{0}\) 处可微且在那里以 \(u\) 为导数的充要条件是有

\[
\lim _ {h \rightarrow 0, h \neq 0} \frac {f (x _ {0} + h) - f (x _ {0}) - u (h)}{\| h \|} = 0.
\]

2.1.2. \(f\) 在 \(x_{0}\) 处严格可微的充要条件是有

\[
\lim_{\substack{(h,k)\to (0,0)\\ h\neq k}}\frac{f(x_{0} + h) - f(x_{0} + k) - \mathrm{D}f(x_{0})\cdot(h - k)}{\|h - k\|} = 0.
\]

2.1.3. 设 \(F_{1}\) 和 \(F_{2}\) 是两个分离局部凸空间，u 是从 \(F_{1} \times F_{2}\) 到 F 的双线性映射，并满足下列连续性条件：

(SC) 若 \(((a_{n}, b_{n}))\) 是 \(F_{1} \times F_{2}\) 中收敛于元素 \((a, b) \in \mathbf{F}_{1} \times \mathbf{F}_{2}\) 的元素序列，则序列 \((u(a_{n}, b_{n}))\) 在 F 中收敛于 \(u(a, b)\)。

设 \(f_{i}\)（\(i = 1, 2\)）是从 E 中一点 \(x_{0}\) 的某邻域到 \(F_{i}\) 的映射。若 \(f_{1}\) 和 \(f_{2}\) 在 \(x_{0}\) 处可微，则 \(u(f_{1}, f_{2})\) 在 \(x_{0}\) 处可微，且有：

\[
\mathbf {D} (u (f _ {1}, f _ {2})) (x _ {0}). h = u (\mathbf {D} f _ {1} (x _ {0}). h, f _ {2} (x _ {0})) + u (f _ {1} (x _ {0}), \mathbf {D} f _ {2} (x _ {0}). h)
\]

对所有 \(h \in E\) 均成立。

\subsection{中值定理}

2.2.1. 设 \(x,y\) 属于 \(E\)，\([x,y]\) 是连接这两点的闭线段。此外，设 \(f\) 是从 \([x,y]\) 的某邻域到空间 \(F\) 的映射，并在 \([x,y]\) 的每一点处可微。则 \(f(x)-f(y)\) 属于当 \(Df(z).(x-y)\) 遍历 \(z\) 时的点集 \([x,y]\) 的闭凸包。

2.2.2. 设 U 是 E 的连通开子集，f 是从 U 到
F 的映射，并在 U 的每一点处导数为零；则 f 在 U 上为常值。

2.2.3. 设 U 是 E 的凸开子集，f 是从 U 到 F 的映射，并在 U 的每一点处可微。给定 F 上的连续半范数 γ 和实数 M ≥ 0，下列条件等价：

\begin{enumerate}
\def\labelenumi{(\roman{enumi})}
\item
  对每个 \(x\)  \(\mathbf{U}\)，都有 \(\| Df(x)\|_{\gamma}\leqslant M\)。
\item
  对 U 中任意 \(x\) 和任意 \(y\)\(U\)，都有 \(\| f(x) - f(y)\|_{\gamma} \leqslant M . \| x - y\| \)。
\end{enumerate}

2.2.4. 设 U 是 E 中一点 \(x_0\) 的一个邻域，\(f\) 是在 U 去掉 \(x_0\) 的补集上定义且取值于 F 的函数。假设 \(f\) 在 U 中每个满足 \(Df(x)\) 的点 \(x\) 处都有导数 \(x \neq x_0\)，且函数 \(x \mapsto Df(x)\) 在 \(D_0\) 趋于 \(x\) 时有极限 \(x_0\)。那么，当  (E) \(\geqslant 2\) 时，\(f\) 在 \(x_0\) 处有极限，并且连续延拓到整个 U 的函数 \(f\) 在 x＿0 处可微、导数为 \(D_0\)；若 \(x_0\) 且假设 \(f\) 在 \(x_0\) 处有极限，同样结论也成立。

\subsection{\texorpdfstring{\(C^{r}\) 类函数 ( \(r \neq \omega\) )}{C\^{}\{r\} 类函数 (r ≠ ω)}}

2.3.1. 设 U 是 E 的开子集，f 是从 U 到 F 的映射。对 \(C^{r}\)，按 r 归纳如下定义关系“f 属于 \(r \in N\) 类”：

\begin{enumerate}
\def\labelenumi{\arabic{enumi})}
\item
  \(f\) 属于 \(\mathbf{C}^{0}\) 类，当且仅当 f 连续；
\item
  若 \(r\) 是整数且 r \(\geqslant 1\)，则函数 \(f\) 属于 \(\mathbf{C}^{r}\) 类，当且仅当它在 \(\mathbf{U}\)\(Df\) 的每一点处可微，且从 \(\mathbf{U}\) 到 \(\mathcal{L}(\mathbf{E};\mathbf{F})\) 的导数 Df 属于 \(\mathbf{C}^{r-1}\) 类。
\end{enumerate}

属于 C \(^{r}\) 类的函数也称为 r 次连续可微函数。

若 \(f\)\(\mathbf{C}^{\infty}\) 对每个整数 r 都属于 \(\mathbf{C}^{r}\) 类，则称 f 属于 \(r\) 类（或称无穷次可微）。

若 \(f\) 在 U 中属于 \(\mathbf{C}^{r}\) 类，则对每个整数 \(\mathbf{U}\)，\(f\) 都 \(p\)\(p \leqslant r\) 次可微，且函数 \(\mathbf{D}^{p}f\) 属于 \(\mathbf{C}^{r-p}\) 类。

2.3.2. 从 E 的开子集 U 到 F 的 \(C^{r}\) 类映射构成所有从 U 到 F 的映射空间的向量子空间 \(\mathcal{C}^{r}(U;F)\)。当 \(\mathcal{C}^{s}(U;F)\subset\mathcal{C}^{r}(U;F)\) 时，有 \(s\geqslant r\)。

2.3.3. 函数 \(f\) 在 E 的开子集 U 中属于 \(\mathbf{C}^{1}\) 类的充要条件是它在 U 的每一点处严格可微。

若 E 是赋范空间 \(E_{i}\) 的乘积，则从 E 的开集 V 到 F 的映射 \(f\) 属于 \(C^{r}\) 类，当且仅当对每个整数 m  r，\(f\) 的迭代偏导数 \(D_{i_{1}}\ldots D_{i_{m}}f\)\(m \leqslant r\) 都存在且连续。

2.3.4. 设 G 是赋范空间，U 是 E 的开子集。设 V 是 G 的开子集，\(g \in \mathcal{C}^{r}(U; G)\) 且 \(f \in \mathcal{C}^{r}(V; F)\)。若 \(g(U) \subset V\)，则从 U 到 F 的映射 \(f \circ g\) 属于 \(C^{r}\) 类。

设 \(\mathbf{F}_1\) 和 \(\mathbf{F}_2\) 是两个分离局部凸空间，\(u\) 是从 \(\mathbf{F}_1 \times \mathbf{F}_2\) 到 \(\mathbf{F}\) 的双线性映射，并且关于 \(\mathbf{F}_1\)（分别地，\(\mathbf{F}_2\)）的有界子集族是半连续的（Topological Vector Spaces，第 III 章，第 4 节，第 2 号）。设 \(\mathbf{U}\) 是 \(\mathbf{E}\) 的开子集，且 \(f_i \in \mathcal{C}^r(\mathbf{U};\mathbf{F}_i)\)（\(i = 1,2\)）。则函数 \(u(f_1,f_2)\) 属于 \(\mathcal{C}^r(\mathbf{U};\mathbf{F})\)。若 \(\mathbf{E}\) 有限维，只需假设 \(u\) 满足第 2.1.3 号的条件 (SC)。

2.3.5. 若 E 是赋范空间 \(E_{i}\)（\(1 \leqslant i \leqslant n\)）的乘积，且 f 是从 E 到 F 的连续 n 线性映射，则 f 属于 \(C^{\infty}\) 类，并且当 \(D^{p}f = 0\) 时有 \(p \geqslant n + 1\)。

2.3.6. 假设 E 和 F 是巴拿赫空间。设 f 是在 E 中一点 \(C^{r}\) 的某邻域上定义且取值于 F 的 \(r \geqslant 1\) 类函数（\(x_{0}\)）。设 \(y_{0} = f(x_{0})\)，并假设 \(Df(x_{0})\) 是从 E 到 F 的同构。则 f 诱导出一个从 \(x_{0}\) 的某邻域到 \(y_{0}\) 的某邻域的同胚 g（1.5），且 g 的逆映射在 \(C^{r}\) 的某邻域上属于 \(y_{0}\) 类。

\subsection{\texorpdfstring{\(C^{r}\) 类函数的导数}{C\^{}\{r\} 类函数的导数}}

2.4.1. 设 \(f\) 是从 E 的开集 U 到 F 的 \(\mathbf{C}^{r}\) 类映射。对每个 \(x \in \mathbf{U}\) 以及每个满足 \(s\)\(s \leqslant r\)  r 的整数 s，多线性映射 \(\mathbf{D}^{s}f(x)\) 都是对称的。

2.4.2. 再设 E 有限维，且 \((e_{1},\ldots,e_{n})\) 是 E 的一组基。偏导数 \(\partial_{i_{1}}\ldots\partial_{i_{s}}f\) 关于指标 \(i_{1},\ldots,i_{s}\) 对称地依赖。设指标 k 在序列 \(\alpha_{k}\) 中出现的次数为 \(i_{1},\ldots,i_{s}\)，并令 \(\alpha=(\alpha_{1},\ldots,\alpha_{n})\)。于是置：

\[
\partial^ {\alpha} f = \partial_ {1} ^ {\alpha_ {1}} \dots \partial_ {n} ^ {\alpha_ {n}} f = \partial_ {i _ {1}} \dots \partial_ {i _ {s}} f
\]

当相对于基 \((e_{1},\ldots,e_{n})\) 的坐标记为 \(x_{1},\ldots,x_{n}\) 时，也将 \(\partial^{\alpha}f\) 写成：

\[
\frac {\partial^ {| \alpha |} f}{\partial x _ {1} ^ {\alpha_ {1}} \dots \partial x _ {n} ^ {\alpha_ {n}}}
\]

\subsection{泰勒公式}

2.5.1. 设 \(r\) 是满足 r \(\geqslant 1\) 的整数，\(f\) 是从 E 的开集 U 到 F 的 \(C^r\) 类映射。对于 \(x \in U\)、\(h \in E\) 和 \(p \leqslant r\)，约定用 \(D^{p}f(x_0).h^p\) 代替 \(D^{p}f(x_0).(h,\ldots,h)\)。若线段 \([x,x + h]\) 包含于 U，则有公式（“泰勒公式”）：

\[
f (x + h) = \sum_ {p = 0} ^ {r - 1} \frac {1}{p !} \mathbf {D} ^ {p} f (x). h ^ {p} + v _ {r} (x; h)
\]

其中“余项” \(v_{r}(x; h)\) 由下式给出：

\[
v _ {r} (x; h) = \int_ {0} ^ {1} \frac {(1 - t) ^ {r - 1}}{(r - 1) !} \mathrm{D} ^ {r} f (x + t h). h ^ {r} d t
\]

有：

\[
v _ {r} (x; h) \equiv \frac {1}{r !} \mathrm{D} ^ {r} f (x). h ^ {r} \quad \text {   mod   } o (\| h \| ^ {r}) \quad \text {   when   } h \text {   tends   to   } 0
\]

以及

\[
f (x + h) \equiv \sum_ {p = 0} ^ {r} \frac {1}{p !} \mathrm{D} ^ {p} f (x). h ^ {p} \quad \mod o (\| h \| ^ {r})
\]

当 h 趋于零时。

2.5.2. 再设 \(\gamma\) 是 \(F\) 上的连续半范数；若对线段 \(\|D^{r}f(z)\|_{\gamma}\leqslant M\) 的每一点 \(z\) 都有 \([x,x + h]\)，则有：

\[
\| v _ {r} (x; h) \| _ {\gamma} \leqslant \frac {\mathbf {M}}{r !} \| h \| ^ {r}
\]

2.5.3. 再假设 \(\mathbf{E} = \mathbf{R}^n\)。于是有：

\[
f (x + h) \equiv \sum_ {| \alpha | \leqslant r} \Delta^ {\alpha} f (x) h ^ {\alpha} \quad \mod o (\| h \| ^ {r}) \text {   when   } h \text {   tends   to   } 0
\]

其中置：

\[
\Delta^ {\alpha} f (x) = \frac {1}{\alpha !} \partial^ {\alpha} f (x)
\]

2.5.4. 设 \(f\) 和 \(g\)\(C^r\) 是 E 的开集 \(U\)\(E\) 上取值于 \(F\) 的两个 C＾r 类函数。\(f\) 和 \(g\) 在 U 中一点 \(x\)\(U\) 处的接触阶 \(\geqslant r\) 的充要条件是，对每个满足 \(D^p f(x) = D^p g(x)\) 的整数 \(p\)，都有 \(0 \leqslant p \leqslant r\)。当 \(E\)\(\leqslant r\) 有限维时，这就是说 \(f\) 和 \(g\) 的  r 阶迭代偏导数（相对于 \(E\) 的一组基）在点 \(x\) 处相等。

2.5.5. 设 U 是 \(E \times R^{n}\) 的一个开子集，形如 \(V \times I_{1} \times \cdots \times I_{n}\)，其中 V 是 E 的开子集，\(I_{1}, \ldots, I_{n}\) 是包含 0 的 R 的开区间。置 \(U_{0} = V\)，并对 \(U_{j} = V \times I_{1} \times \cdots \times I_{j}\) 置 \(1 \leqslant j \leqslant n\)。给定函数 \(f \in \mathcal{C}^{r}(U; F)\)（\(1 \leqslant r \leqslant \infty\)），存在唯一的函数序列 \(f_{j}\in\mathcal{C}^{r-1}(\mathbf{U}_{j};\mathbf{F})\)（\(0\leqslant j\leqslant n\)），使得：

\[
f (x, t _ {1}, \dots , t _ {n}) = f _ {0} (x) + \sum_ {j = 1} ^ {n} t _ {j} f _ {j} (x, t _ {1}, \dots , t _ {j})
\]

对 \(x \in V\) 和 \(t_j \in I_j\) 成立。有：

\[
f _ {0} (x) = f (x, 0, \dots , 0)
\]

\[
f _ {j} (x, t _ {1}, \dots , t _ {j}) = \int_ {0} ^ {1} \partial_ {j} f (x, t _ {1}, \dots , t _ {j - 1}, t _ {j} u, 0, \dots , 0) d u
\]

对 \(1 \leqslant j \leqslant n\) 成立。在最后一个公式中，\(\partial_j f\) 表示函数 (t＿1, , t＿n)  f(x, t＿1, , t＿n) 的第 \(j\)\((t_1, \ldots, t_n) \mapsto f(x, t_1, \ldots, t_n)\) 个偏导数。

\subsection{可微性判据}

2.6.1. 假设 F 除了拓扑 T 外还赋有较粗拓扑 T'，且 T' 也使 F 成为分离局部凸空间。再假设 T 和 T' 满足下列条件：

\begin{enumerate}
\def\labelenumi{(\Alph{enumi})}
\setcounter{enumi}{18}
\tightlist
\item
  对 T 拓扑下 0 的每个邻域 V，都存在 T 拓扑下 0 的一个邻域 W，使得 W 的每个 T-紧子集的 T'-闭凸包都包含于 V。
\end{enumerate}

设 \(f\) 是从 E 的开子集 U 到 F 的映射。假设 F 赋予拓扑 T' 时 \(f\) 属于 \(\mathbf{C}^{r}\) 类（\(1 \leqslant r \leqslant \infty\)）；对 U 中所有 x 以及每个整数 \(\mathcal{T}'\)，\(\mathrm{D}^{m}f(x)\) 都是从 \(x\) 到赋予拓扑 \(m \leqslant r\) 的 F 的连续多线性映射；且映射 \(\mathrm{E}^{m}\) 从 U 到 \(\mathcal{T}\) 连续（F 赋予拓扑 \(x \mapsto \mathrm{D}^{m}f(x)\)\(\mathcal{L}_{m}(\mathrm{E};\mathrm{F})\)）。则 F 赋予拓扑 \(\mathcal{T}\) 时 \(f\) 属于 \(\mathbf{C}^{r}\)\(\mathcal{T}\) 类，并且其导数 \(\mathrm{D}^{m}f\) 对 \(\mathcal{T}\) 和 \(\mathcal{T}'\) 相同。

条件 (S) 特别地在如下情形成立：存在 T 拓扑下 0 的邻域基本系，其成员对拓扑 T' 都是闭的；当 F 赋予 T 时的对偶与 F 赋予 T' 时的对偶相同时，这一情形确实发生。若 F 赋予拓扑 T 时拟完备，条件 (S) 也成立（Topological Vector Spaces，第 III 章，第 2 节，第 5 号）。

2.6.2. 设 \(f\) 是从 E 的开集 \(U\)\(E\) 到 \(F\) 的映射。若 \(f\) 属于 \(C^r\) 类（\(0 \leqslant r \leqslant \infty\)），则对 F 上每个连续线性形式 u，标量函数 \(u \circ f\) 都属于 \(C^r\) 类。反之，若 F 拟完备，且对所有 \(u\)  \(F\)'，函数 \(u \circ f\) 都属于 \(C^{r+1}\)\(u \in F'\) 类，则 \(f\) 属于 \(C^r\) 类。

\section{实解析或复解析函数}

本段假定 K = R 或 C。用 \((\mathrm{E}_{i})\)（\(1 \leqslant i \leqslant n\)）表示 K 上赋范空间的一个有限族，用 E 表示各 \(E_{i}\) 的乘积空间，用 F 表示 K 上的分离局部凸拓扑向量空间。

\subsection{收敛级数}

3.1.1. 设 \(f = \sum f_{\alpha}\) 是属于 \(\hat{\mathbf{P}}(\mathbf{E}_1, \ldots, \mathbf{E}_n; \mathbf{F})\) 的形式级数（附录，第 A.5 号）。若 \(\gamma\) 是 \(\mathbf{F}\) 上的连续半范数，且 \(\mathbf{R} = (\mathbf{R}_1, \ldots, \mathbf{R}_n)\) 是 \(n\) 个严格正实数组成的序列，置：

\[
\| f \| _ {\gamma , \mathbf {R}} = \sum_ {\alpha} \mathbf {R} ^ {\alpha} \| f _ {\alpha} \| _ {\gamma}.
\]

若 \(F\) 是赋范空间且 \(\gamma\) 是 \(F\) 的范数，则用 \(\| f \|_{\mathbb{R}}\) 代替 \(\| f \|_{\gamma, \mathbb{R}}\)。

由满足对 \(\mathcal{H}_{\mathbb{R}}(\mathbf{E}_{1},\ldots,\mathbf{E}_{n};\mathbf{F})\) 上每个连续半范数 \(f\in\hat{\mathbf{P}}(\mathbf{E}_{1},\ldots,\mathbf{E}_{n};\mathbf{F})\)，\(\|f\|_{\gamma,\mathbb{R}}\) 都有限的 \(\gamma\) 构成的集合 \(\mathbf{F}\)，是 \(\hat{\mathbf{P}}(\mathbf{E}_{1},\ldots,\mathbf{E}_{n};\mathbf{F})\) 的向量子空间。有

\[
\mathcal {H} _ {\mathbb {R}} (\mathrm{E} _ {1}, \dots , \mathrm{E} _ {n}; \mathrm{F}) \subset \mathcal {H} _ {\mathbb {R} ^ {\prime}} (\mathrm{E} _ {1}, \dots , \mathrm{E} _ {n}; \mathrm{F})
\]

只要 \(R_{i} \geqslant R_{i}'\)（\(1 \leqslant i \leqslant n\)）即可。下列各空间的并

\[
\mathcal {H} _ {\mathbf {R}} (\mathrm{E} _ {1}, \dots , \mathrm{E} _ {n}; \mathrm{F})
\]

\(\mathcal{H}(E_{1},\ldots,E_{n};F)\)是 \(\hat{\mathrm{P}}(E_{1},\ldots,E_{n};F)\) 的向量子空间，记作 \(E_{i}\)，其中元素称为取值于 \(F\) 的各 \(E_{i}\) 的乘积上的收敛级数。它只取决于各 \(\mathcal{H}(E_{1},\ldots,E_{n};F)\) 的拓扑而不取决于它们的范数。因此，当各 \(E_{i}\) 是可赋范空间时，无须在每个 \(E_{i}\) 上选取范数，也可以谈论空间 H(E＿1,,E＿n;F)。

映射 \(f\mapsto\|f\|_{\gamma,\mathbb{R}}\) 是 \(\mathcal{H}_{\mathbb{R}}(\mathrm{E}_{1},\ldots,\mathrm{E}_{n};\mathrm{F})\) 上的半范数。当 \(\gamma\) 遍历 F 上的连续半范数（或仅遍历定义 F 的拓扑的一组半范数）时，由这些半范数定义的拓扑是豪斯多夫的。若 F 可赋范（分别地，完备），则 \(\mathcal{H}_{\mathbb{R}}(\mathrm{E}_{1},\ldots,\mathrm{E}_{n};\mathrm{F})\) 也可赋范（分别地，完备）。\(\mathcal{H}(\mathrm{E}_{1},\ldots,\mathrm{E}_{n};\mathrm{F})\) 上那些限制到每个 \(\mathcal{H}_{\mathbb{R}}\) 都连续的半范数，在 \(\mathcal{H}(\mathrm{E}_{1},\ldots,\mathrm{E}_{n};\mathrm{F})\) 上定义一个豪斯多夫拓扑，该拓扑只取决于各 \(\mathrm{E}_{i}\) 的拓扑。从 \(\mathcal{H}(\mathrm{E}_{1},\ldots,\mathrm{E}_{n};\mathrm{F})\) 到 \(\hat{\mathrm{P}}(\mathrm{E}_{1},\ldots,\mathrm{E}_{n};\mathrm{F})\) 的嵌入是连续的。

3.1.3. \(\hat{\mathbf{P}}(\mathbf{E};\mathbf{F})\) 到 \(\hat{\mathbf{P}}(\mathbf{E}_{1},\ldots,\mathbf{E}_{n};\mathbf{F})\) 的规范同构通过限制给出从 \(\mathcal{H}(\mathbf{E};\mathbf{F})\) 到 \(\mathcal{H}(\mathbf{E}_{1},\ldots,\mathbf{E}_{n};\mathbf{F})\) 的拓扑向量空间同构。

3.1.4. 设 \(f \in \mathcal{H}(\mathrm{E}_1, \ldots, \mathrm{E}_n; \mathrm{F})\)，并令 \(J(f)\)\(R \in (\mathbf{R}_+^*)^n\) 为满足 \(f \in \mathcal{H}_{\mathbb{R}}(\mathrm{E}_1, \ldots, \mathrm{E}_n; \mathrm{F})\) 的 \(I(f)\) 的集合。\(J(f)\) 的内部 I(\(f\)) 称为 f 的严格收敛指标。它由满足存在 \(R\)\(R' \in J(f)\)'  J(f) 且 \(0 < R_i < R'_i\)（\(1 \leqslant i \leqslant n\)）的 R 构成。用 \(\Omega(f)\)\((\log R_1, \ldots, \log R_n)\) 表示 R 属于 I(f) 时 \(R^n\) 中的点 \(R \in I(f)\) 的集合：它是 \(R^n\) 的凸子集。

当 \(n=1\) 时，集合 \(\mathbf{I}(f)\) 是 \([0,\rho(f)]\) 中的区间 \(\mathbf{R}\)，\(\rho(f)\) 称为 \(f\)\(+\infty\) 的严格收敛半径。它也是满足如下条件的实数 \(R>0\) 的集合的上确界（有限或 \(\gamma\)）：对 \(F\) 上每个连续半范数 ，存在常数 \(M\)，使得（其中 \(\|f_{m}\|_{\gamma}\leqslant MR^{-m}\)）对每个整数 \(f=\sum f_{m},f_{m}\in P_{m}(E;F)\) 都有 \(m\geqslant0\)。

满足存在 \(x = (x_i) \in \mathbf{E}_1 \times \cdots \times \mathbf{E}_n\) 且对所有 i 有 \(\mathbf{R} \in \mathbf{I}(f)\) 的点 \(\|x_i\| \leqslant R_i\) 的集合称为 f 的严格收敛域，记作 \(\mathbf{C}(f)\)。它也是满足存在 \(\mathbf{R} \in \mathbf{J}(f)\) 且对所有 i 有 \(\|x_i\| < R_i\) 的点 x 的集合的内部。

3.1.5. 对于 \(f_{\alpha} \in \mathbf{P}_{\alpha}(\mathbf{E}_1, \ldots, \mathbf{E}_n; \mathbf{F})\) 以及 \(\gamma\) 上每个连续半范数 \(\mathbf{F}\)，置：

\[
\| f _ {\alpha} \| _ {\gamma} = \sup _ {\| x _ {i} \| \leqslant 1} \| f _ {\alpha} (x _ {1}, \dots , x _ {n}) \| _ {\gamma}.
\]

于是有不等式：

\[
\left\| \left| f _ {\alpha} \right| \right\| _ {\gamma} \leqslant \left\| f _ {\alpha} \right\| _ {\gamma} \leqslant \frac {\alpha^ {\alpha}}{\alpha !} \left\| \left| f _ {\alpha} \right| \right\| _ {\gamma}.
\]

对于 \(f = \sum f_{\alpha} \in \hat{\mathbf{P}}(\mathrm{E}_1, \ldots, \mathrm{E}_n; \mathbf{F})\) 和 \(\mathbf{R} \in (\mathbf{R}_{+}^{*})^n\)，置：

\[
\| | f \| | _ {\gamma , \mathbf {R}} = \sum_ {\alpha} \mathbf {R} ^ {\alpha} \| | f _ {\alpha} \| | _ {\gamma}
\]

并用 \(\tilde{\mathcal{H}}_{\mathbb{R}}(\mathrm{E}_{1},\ldots ,\mathrm{E}_{n};\mathrm{F})\) 表示下列空间的向量子空间

\[
\hat {\mathbf {P}} (\mathbf {E} _ {1}, \dots , \mathbf {E} _ {n}; \mathbf {F})
\]

由满足对每个  都有  | \(f\)\(\| | f \| | _ {\gamma,\mathbb{R}}\)  | ＿ \(\gamma\),R 有限的 f 构成，并赋予由半范数 \(f\mapsto\| | f \| | _ {\gamma,\mathbb{R}}\) 定义的拓扑。有 \(\mathcal{H}_{\mathbb{R}}\subset\tilde{\mathcal{H}}_{\mathbb{R}}\)，且从 \(\mathcal{H}_{\mathbb{R}}\) 到 \(\tilde{\mathcal{H}}_{\mathbb{R}}\) 的嵌入是连续的。空间 \(\mathcal{H}(\mathrm{E}_{1},\ldots,\mathrm{E}_{n};\mathrm{F})\) 是各 \(\tilde{\mathcal{H}}_{\mathbb{R}}(\mathrm{E}_{1},\ldots,\mathrm{E}_{n};\mathrm{F})\) 的并，其拓扑是使 \(\mathcal{H}_{\mathbb{R}}\) 到 \(\mathcal{H}\) 的嵌入连续的最细局部凸拓扑。

若 \(f \in \mathcal{H}(E_1, \ldots, E_n; F)\)，则称满足 \(\tilde{I}(f)\) 的 \(\tilde{J}(f)\) 的集合 \(R \in (\mathbf{R}_+^*)\) 的内部 \(f \in \mathcal{H}_{\mathbb{R}}\) 为收敛指标。有：

\[
e ^ {- 1} \mathbf {I} (f) \subset \tilde{\mathbf {I}} (f) \subset \mathbf {I} (f)
\]

（其中 e 是自然对数的底）。从 \(\tilde{I}(f)\) 出发，按 3.1.4 定义收敛域 \(\tilde{C}(f)\)，并在 \(n = 1\) 时定义收敛半径 \(\tilde{\rho}(f)\)。特别地，有：

\[
e ^ {- 1} \tilde {\rho} (f) \leqslant \rho (f) \leqslant \tilde {\rho} (f).
\]

3.1.6. 对于 \(\mathbf{R} \in (\mathbf{R}_{+}^{*})^{n}\)，将 E 中以 0 为中心、半径为 \(\mathbf{R}\) 的（闭）多球称为由满足对所有 i 都有 \(\mathbf{B}(\mathbf{R})\) 的 \(x \in \mathbf{E}\) 构成的集合 \(\|x_{i}\| \leqslant \mathbf{R}_{i}\)。若  E＿\(i\)\(\dim \mathbf{E}_{i} = 1\) = 1，也称为多圆盘。若 \(f \in \mathcal{H}(\mathbf{E}_{1}, \ldots, \mathbf{E}_{n}; \mathbf{F})\)，则 \(f\) 的收敛域（分别地，严格收敛域）是 \(\mathbf{B}(\mathbf{R})\)（分别地，\(\mathbf{R} \in \tilde{\mathbf{I}}(f)\)）时多球 \(\mathbf{R} \in \mathbf{I}(f)\) 的并。

3.1.7. 设 \(f \in \mathcal{H}(E_1, \ldots, E_n; F)\)。假设 F 拟完备。对每个 \(x \in \tilde{\mathbf{C}}(f)\)，族 \(f_{\alpha}(x)\) 在 F 中可求和。其和记作 \(\hat{f}(x)\)，或简写为 \(f(x)\)，是 \(\tilde{\mathbf{C}}(f)\) 上的连续函数。更确切地说，对每个满足 \(f \in \tilde{\mathcal{H}}_{\mathbb{R}}\) 的 R，族 \(f_{\alpha}(x)\) 当 \(x \in B(R)\) 时一致可求和。映射 \(f \mapsto \hat{f}\) 是从 \(\tilde{\mathcal{H}}_{\mathbb{R}}\) 到 \(B(R)\) 上有界连续函数空间的连续线性单射，后者赋予一致收敛拓扑。

3.1.8. 设 \(F_1, \ldots, F_m\) 是分离多范数空间，\(u\)\(m\) 是从 \(F_1 \times \cdots \times F_m\) 到 \(F\) 的连续 m 线性映射。设 \(f_i \in \mathcal{H}(E_1, \ldots, E_n; F_i)\)（\(1 \leqslant i \leqslant m\)）。形式级数 \(u(f_1, \ldots, f_m)\) 属于 \(\mathcal{H}(E_1, \ldots, E_n; F)\)。有：

\[
\mathrm{C} (u (f _ {1}, \dots , f _ {m})) \supset \bigcap_ {i} \mathrm{C} (f _ {i})
\]

\[
\tilde {C} (u (f _ {1}, \dots , f _ {m})) \supset \bigcap_ {i} \tilde {C} (f _ {i})
\]

并且，若 \(x \in \bigcap_{i} \tilde{C}(f_i)\)，且 F 和各 \(F_i\) 拟完备，则有：

\[
u (f _ {1}, \dots , f _ {m}) (x) = u (f _ {1} (x), \dots , f _ {m} (x)).
\]

3.1.9. 设 \(\mathbf{F}_1, \ldots, \mathbf{F}_m\) 是完备赋范空间，且 \(\mathbf{F}\) 拟完备。令 \(\mathbf{f} = (f_i)_{1 \leqslant i \leqslant m}\)，其中 \(f_i \in \mathcal{H}(\mathbf{E}_1, \ldots, \mathbf{E}_n; \mathbf{F}_i)\)，\(g \in \mathcal{H}(\mathbf{F}_1, \ldots, \mathbf{F}_m; \mathbf{F})\)，并假设 \((f_i(0))_{1 \leqslant i \leqslant m}\) 属于 \(g\) 的严格收敛域。于是，对每个 \(\alpha \in \mathbb{N}^m\)，形式级数 \(g_\alpha \circ \mathbf{f}\) 属于 \(\mathcal{H}(\mathbf{E}_1, \ldots, \mathbf{E}_n; \mathbf{F})\)，且族 \(g_\alpha \circ \mathbf{f}\) 在 \(\mathcal{H}(\mathbf{E}_1, \ldots, \mathbf{E}_n; \mathbf{F})\) 中可求和，从而当然也在 \(\hat{\mathbf{P}}(\mathbf{E}_1, \ldots, \mathbf{E}_n; \mathbf{F})\) 中可求和。其和记作 \(g \circ \mathbf{f}\)。

更确切地说，存在 \(\mathbb{R}\in\bigcap_{i}\mathrm{I}(f_{i})\) 和 \(\mathbb{R}^{\prime}\in\mathrm{I}(g)\)，使得当 \(\|f_{i}\|_{\mathbb{R}}<R_{i}^{\prime}\) 时 \(1\leqslant i\leqslant m\)。在这些条件下，形式级数 \(g_{\alpha}\circ f\) 属于 \(\mathcal{H}_{\mathbb{R}}(\mathrm{E}_{1},\ldots,\mathrm{E}_{n};\mathrm{F})\)，且族 \(g_{\alpha}\circ f\) 在 \(\mathcal{H}_{\mathbb{R}}(\mathrm{E}_{1},\ldots,\mathrm{E}_{n};\mathrm{F})\) 中可求和。最后，若 \(x\in\mathrm{B}(\mathbb{R})\)，则 \(\mathbf{f}(x)=(f_{i}(x))\) 属于 \(\mathrm{B}(\mathrm{R}^{\prime})\subset\mathrm{F}_{1}\times\cdots\times\mathrm{F}_{m}\)，且有：

\[
g (\mathbf {f} (x)) = (g \circ \mathbf {f}) (x).
\]

3.1.10. 假设 \(E_{i} = K\)（\(1 \leqslant i \leqslant n\)）。于是将空间 \(\hat{\mathbf{P}}(\mathbf{K}^{n}; \mathbf{F})\) 识别为具有 F 中系数的 n 个不定元 \(X_{1}, \ldots, X_{n}\) 的形式幂级数空间，且 \(\hat{\mathbf{P}}(\mathbf{K}^{n};\mathbf{F})\) 中的元素写作：

\[
f = \sum_ {\alpha} X ^ {\alpha} c _ {\alpha} \quad \text {   with   } c _ {\alpha} \in F.
\]

若 \(\mathbf{R} \in (\mathbf{R}_{+}^{*})^{n}\)，且 \(\gamma\) 是 \(\mathbf{F}\) 上的连续半范数，则有：

\[
\| | f \| | _ {\gamma , \mathbf {R}} = \| f \| _ {\gamma , \mathbf {R}} = \sum_ {\alpha} \mathbf {R} ^ {\alpha} \| c _ {\alpha} \| _ {\gamma}
\]

有 \(\mathbf{I}(f) = \tilde{\mathbf{I}}(f)\)、\(\mathbf{C}(f) = \tilde{\mathbf{C}}(f)\)，并且当 \(n = 1\) 时 \(\rho(f) = \tilde{\rho}(f)\)。

系数在 \(\mathcal{H}(\mathbf{K}^{n};\mathbf{K})\) 中的收敛级数空间 \(\mathbf{K}\) 也记作 \(\mathbf{K}\{\{\mathbf{X}_{1},\ldots ,\mathbf{X}_{n}\}\}\)；它是 \(\mathbf{K}[[\mathbf{X}_1,\dots ,\mathbf{X}_n]]\) 的子代数。空间 \(\mathcal{H}(\mathbf{K}^n;\mathbf{F})\) 是 \(\mathbf{K}\{\{\mathbf{X}_1,\ldots ,\mathbf{X}_n\}\}\) 上的模；若 \(\mathbf{F}\) 有限维，则将此模识别为 \(\mathbf{K}\{\{\mathbf{X}_1,\ldots ,\mathbf{X}_n\}\} \otimes_{\mathbf{K}}\mathbf{F}\)。

3.1.11. 设 \(\mathbf{f} \in \mathcal{H}(\mathbf{K}^n; \mathbf{K}^m)\)，它由系数在 K 中的 \(m\) 个收敛级数 \(f_j(X_1, \ldots, X_n)\) 组成的系统表示。同样，设 \(\mathbf{g} \in \mathcal{H}(\mathbf{K}^m; \mathbf{K}^p)\)，它由 \(p\) 个收敛级数 \(g_k(Y_1, \ldots, Y_m) = \sum g_{k,\beta} Y^\beta\) 组成的系统表示。空间 \(\mathbf{h} = \mathbf{g} \circ \mathbf{f}\) 中的元素 \(\mathcal{H}(\mathbf{K}^n; \mathbf{K}^p)\)（参见 3.1.9）由 \(p\) 个形式级数 \(h_k(X_1, \ldots, X_n)\) 表示，其确定如下：对于 \(\alpha \in \mathbf{N}^n\) 和 \(\beta \in \mathbf{N}^m\)，令 \(c_{\alpha,\beta}\) 为形式级数 \(X^\alpha\) 中 \(\mathbf{f}^\beta = \prod f_j^{\beta_j}\) 的系数；则族 \((g_{k,\beta} c_{\alpha,\beta})_{\beta \in \mathbb{N}^m}\) 在 K 中可求和，其和就是 \(X^\alpha\) 中 \(h_k\) 的系数。

3.1.12. 假设 F 拟完备，且 \(\hat{\mathbf{E}}_{i}\) 是 \(\mathbf{E}_{i}\) 的完备化。在 \(\mathbf{E}_{1} \times \cdots \times \mathbf{E}_{n}\) 上取值于 F 的连续多项式可由连续性延拓为在 \(\hat{\mathbf{E}}_{1} \times \cdots \times \hat{\mathbf{E}}_{n}\) 上取值于 F 的多项式连续映射。由此得到从 \(j\) 到 \(\hat{\mathbf{P}}(\mathbf{E}_{1}, \ldots, \mathbf{E}_{n}; \mathbf{F})\) 的双射 \(\hat{\mathbf{P}}(\hat{\mathbf{E}}_{1}, \ldots, \hat{\mathbf{E}}_{n}; \mathbf{F})\)。若 \(f \in \mathcal{H}(\mathbf{E}_{1}, \ldots, \mathbf{E}_{n}; \mathbf{F})\)，则 \(j(f) \in \mathcal{H}(\hat{\mathbf{E}}_{1}, \ldots, \hat{\mathbf{E}}_{n}; \mathbf{F})\)，反之亦然。\(f\) 和 \(j(f)\) 的严格收敛指标相同。

3.1.13. 假设 K = R，但 F 赋有与其 R 向量空间结构相容的复向量空间结构。令 \(E_{i}^{C} = E_{i} \otimes_{R} C\)。若 \(y \in E_{i}^{C}\)，置：

\[
\| y \| = \inf \sum_ {k} | a _ {k} |. \| x _ {k} \|,
\]

下确界取遍满足 \((x_{k}, a_{k}) \in \mathrm{E}_{i} \times \mathbf{C}\) 的所有有限对族 \(y = \sum_{k} x_{k} \otimes a_{k}\)。于是得到复向量空间 \(E_{i}^{c}\) 上的一个范数；若同意将 x  \(E_{i}\) 与 \(x \in E_{i}\) 识别，则此范数延拓给定的 \(x \otimes 1\) 上的范数。设 \(h\) 是 \(E_{1} \times \cdots \times E_{n}\) 上取值于 F、关于多重次数 \(\alpha\) 齐次的 R-多项式连续映射；则存在唯一的在 \(\tilde{h}\) 上取值于 F 的 C-多项式连续映射 \(E_{1}^{c} \times \cdots \times E_{n}^{c}\)，延拓 h 且关于多重次数 \(\alpha\) 齐次。有

\[
\| \tilde {h} \| _ {\gamma} = \| h \| _ {\gamma}
\]

对复向量空间 F 上的每个连续半范数 γ。

\[
f = \sum_ {\alpha} f _ {\alpha} \in \mathcal {H} (\mathrm{E} _ {1}, \dots , \mathrm{E} _ {n}; \mathrm{F}), \text {   then   } \tilde {f} = \sum_ {\alpha} \tilde {f} _ {\alpha} \in \mathcal {H} (\mathrm{E} _ {1} ^ {\mathrm{c}}, \dots , \mathrm{E} _ {n} ^ {\mathrm{c}}; \mathrm{F}).
\]

级数 \(f\) 和 \(\tilde{f}\) 具有相同的严格收敛指标（且当 \(n = 1\) 时具有相同的严格收敛半径）。

反之，假设 K = C。令 \(E_{i}^{0}\) 和 \(F^{0}\) 为通过限制标量得到的 R 上的空间。若 \(f_{\alpha} \in \mathrm{P}_{\alpha}(\mathrm{E}_{1}, \ldots, \mathrm{E}_{n}; \mathrm{F})\)，则 \(f_{\alpha} \in \mathrm{P}_{\alpha}(\mathrm{E}_{1}^{0}, \ldots, \mathrm{E}_{n}^{0}; \mathrm{F}^{0})\)。若 \(f = \sum_{\alpha} f_{\alpha} \in \mathcal{H}(\mathrm{E}_{1}, \ldots, \mathrm{E}_{n}; \mathrm{F})\)，则形式级数 \(f^{0} = \sum_{\alpha} f_{\alpha} \in \hat{\mathrm{P}}(\mathrm{E}_{1}^{0}, \ldots, \mathrm{E}_{n}^{0}; \mathrm{F}^{0})\) 是收敛级数。f 和 \(f^{0}\) 的收敛指标（分别地，严格收敛指标）相同，并且对所有 \(f(x) = f^{0}(x)\)，有 \(x \in \tilde{\mathbf{C}}(f) = \tilde{\mathbf{C}}(f^{0})\)。

\subsection{解析函数}

3.2.1. 设 U 是 E 的开子集，f 是从 U 到 F 的映射。若对 U 中每一点 a，都存在收敛级数 \(C^{\omega}\)，使得当 E 中的 x 充分接近零时 \(f_{a} \in \mathcal{H}(E; F)\)，则称 f 在 U 中属于 \(f(a + x) = f_{a}(x)\) 类，或称 K-解析（或简写为解析）。当 K = R（分别地，C）时，也称 f 为实解析（分别地，复解析或全纯）。从 U 到 F 的解析映射构成所有从 U 到 F 的映射空间的向量子空间，记作 \(\mathcal{C}^{\omega}(U; F)\)。

对于 \(a \in \mathbf{U}\)，形式级数 \(f_{a}\) 是唯一的：称其为 \(f\) 在点 \(a\) 处的幂级数展开。若 \(f_{a} = \sum_{\alpha}(f_{a})_{\alpha}\)（其中 \((f_{a})_{\alpha} \in \mathrm{P}_{\alpha}(\mathrm{E}_{1}, \ldots, \mathrm{E}_{n}; \mathrm{F})\)），置：

\[
\Delta^ {\alpha} f (a) = \left(f _ {a}\right) _ {\alpha}.
\]

3.2.2. 若 \(f \in \mathcal{C}^{\omega}(\mathrm{U};\mathrm{F})\)，则映射 \(\Delta^{\alpha}f: a \mapsto \Delta^{\alpha}f(a)\) 从 U 到

\[
\mathbf {P} _ {\alpha} (\mathrm{E} _ {1}, \dots , \mathrm{E} _ {n}; \mathrm{F})
\]

是解析的。映射 \(\Delta^{\alpha}: f \mapsto \Delta^{\alpha}f\) 是从 \(\mathcal{C}^{\omega}(\mathrm{U};\mathrm{F})\) 到 \(\mathcal{C}^{\omega}(\mathrm{U};\mathrm{P}_{\alpha}(\mathrm{E}_{1},\ldots,\mathrm{E}_{n};\mathrm{F}))\) 的 K-线性映射。因此，对于 \(a \in \mathrm{U}\) 以及 \(\alpha, \beta \in \mathbf{N}^{n}\)，有 \(\Delta^{\beta}(\Delta^{\alpha}f)(a) \in \mathrm{P}_{\beta}(\mathrm{E}_{1},\ldots,\mathrm{E}_{n};\mathrm{P}_{\alpha}(\mathrm{E}_{1},\ldots,\mathrm{E}_{n};\mathrm{F}))\)。若 \(x = (x_{i}) \in \mathrm{E}\)，则 \((\Delta^{\beta}(\Delta^{\alpha}f)(a))(x) \in \mathrm{P}_{\alpha}(\mathrm{E}_{1},\ldots,\mathrm{E}_{n};\mathrm{F})\)，且 \(((\Delta^{\beta}(\Delta^{\alpha}f)(a))(x))(x) \in \mathrm{F}\)。F 中此元素等于 \(((\alpha, \beta)) (\Delta^{\alpha + \beta}f(a))(x)\)。这表示为：

\[
\Delta^ {\beta} \circ \Delta^ {\alpha} = ((\alpha , \beta)) \Delta^ {\alpha + \beta}.
\]

\(n=1\)3.2.3. \(f\) 和 \(\Delta^{\alpha}f\) 在 U 中同一点 \(a\) 处的幂级数展开的严格收敛指标（以及当 n=1 时的严格收敛半径）相同。

3.2.4. 设 \(f \in \mathcal{C}^{\omega}(\mathrm{U};\mathrm{F})\)。则 \(f\) 在 \(\mathbf{U}\) 中严格可微且无穷次可微（若 \(\mathbf{K} = \mathbf{R}\)，则 \(f\) 在 \(\mathbf{C}^{\infty}\) 中属于 \(\mathbf{U} \) 类）。\(f\) 的迭代导数是解析的，且它们在点 \(a\) 处的值是对称多线性映射。于是可如第 2.4.2 号引入迭代偏导数的记号 \(\mathbf{D}^{a}f\)。有：

\[
\alpha ! \Delta^ {\alpha} f (a) (h) = D ^ {\alpha} f (a). (h, \dots , h)
\]

对所有 \(a \in \mathbf{U}\) 和 \(h \in \mathbf{E}\) 均成立，记作：

\[
\alpha ! \Delta^ {\alpha} f = D ^ {\alpha} f.
\]

3.2.5. 设 U 是 E 的开子集，f、g 是从 U 到 F 的两个解析映射。设 \(a \in U\)。f 和 g 在 a 处接触阶 \(\geqslant k\) 的充要条件是，对所有满足 \(\Delta^{\alpha}f(a) = \Delta^{\alpha}g(a)\) 的 \(\alpha\)，都有 \(|\alpha| \leqslant k\)。若 f 和 g 在 a 处具有无穷阶接触，则它们在 a 的某邻域上相等。U 中 f 和 g 具有无穷阶接触的点集是开且闭的。

特别地，若 U 连通，且存在 U 的一个非空开子集，在其上 f 和 g 相等，则 f = g（“解析延拓原理”）。

3.2.6. 设 U 是 E 的开子集，f 是从 U 到 F 的解析映射。若 f 的导数 Df 为零，则 f 局部为常值。

3.2.7. 假设 F 拟完备，且 G 是完备赋范空间。设 g 是从 E 的开子集 U 到 G 的解析映射，f 是从 G 的包含 g(U) 的开子集 V 到 F 的解析映射。复合映射 \(f \circ g\) 在 U 中解析。再设 \(0 \in U\) 且 \(g(0) = 0\)。将 f 在 0 处的幂级数展开中的变量替换为 g 在 0 处的幂级数展开，即得到 \(f \circ g\) 在 0 处的幂级数展开（3.1.9）。

3.2.8. 设 \(F_{1},\ldots,F_{m}\) 是分离多范数空间，u 是从 \(F_{1}\times\cdots\times F_{m}\) 到 F 的连续多线性映射。设 U 是 E 的开子集，且 \(f_{i}\in\mathcal{C}^{\omega}(U;F_{i})\)。函数 \(u(f_{1},\ldots,f_{m})\) 是解析的，且其在点 \(a\in U\) 处的幂级数展开是级数 \(u((f_{1})_{a},\ldots,(f_{m})_{a})\)（3.1.8）。

3.2.9. 假设 \(F\) 拟完备。设 \(f \in \mathcal{H}(E_1, \ldots, E_n; F)\)；函数 \(x \mapsto f(x)\)（3.1.7）在开集 \(C(f)\) 中解析，其中 C(\(f\)) 是 f 的严格收敛域。若 \(n = 1\) 且 \(\|a\| < \rho(f)\)，则 \(f\) 在 \(a\) 处的幂级数展开的严格收敛半径至少等于 \(\rho(f) - \|a\|\)。若 \(\rho(f) = +\infty\)，则称 \(f\) 为整函数。

3.2.10. 保持 3.2.9 的假设。若 K = C，则将 C(f) 换为 \(\tilde{C}(f)\)，并将 \(\rho(f)\)（当 \(\tilde{\rho}(f)\) 时）换为 \(n = 1\) 后，3.2.9 的结论仍完全成立。若 K = R，则函数 \(x \mapsto f(x)\) 在 \(\tilde{C}(f)\) 中解析。

3.2.11. 假设 \(E_{i} = K\)（\(1 \leqslant i \leqslant n\)）。设 U 是 \(K^{n}\) 的开子集，且 \(f \in \mathcal{C}^{\omega}(K^{n}; F)\)。若 \(0 \in U\)，且 \(f_{0} = \sum_{\alpha} X^{\alpha} c_{\alpha}\) 是 f 在 0 处的幂级数展开，则在将 \(\Delta^{\alpha} f\) 识别为 F 后，\(P_{\alpha}(K^{n}; F)\) 在 0 处的幂级数展开写作：

\[
(\Delta^ {\alpha} f) _ {0} = \sum_ {\beta} ((\alpha , \beta)) X ^ {\beta} c _ {\alpha + \beta}.
\]

特别地，对 \(1 \leqslant i \leqslant n\) 以及 \(x \in \mathbf{C}(f_0)\)：

\[
\partial_ {i} f (x) = \sum_ {\alpha} (\alpha_ {i} + 1) x ^ {\alpha} c _ {\alpha + \varepsilon_ {i}}.
\]

\subsection{全纯函数}

本节假定 K = C。

3.3.1. 假设 F 拟完备。设 U 是 E 的开子集，f 是从 U 到 F 的映射。下列条件等价：

\begin{enumerate}
\def\labelenumi{(\roman{enumi})}
\item
  \(f\) 是全纯的；
\item
  \(f\) 是可微的；
\item
  \(f\) 局部有界，并且对任意 \(a \in \mathbf{U}\) 和 \(h \in \mathbf{E}\)，在 \(t \mapsto f(a + th)\) 中 0 的某开邻域上定义的函数 \(\mathbf{C}\) 是全纯的；
\item
  \(f\) 局部有界，并且对 F 上每个连续线性形式 \(u\)，取值于 C 的函数 \(u \circ f\) 是全纯的；
\item
  \(f\) 连续且局部有界，并且 F 的对偶中存在一个完备集 H，使得对每个 \(u \circ f\)，\(u \in \mathbf{H}\) 都是全纯的。
\end{enumerate}

当 E 有限维时（分别地，当 F 是巴拿赫空间时），去掉条件 (iii)、(iv) 或 (v)（分别地，条件 (iv) 或 (v)）中的假设“f 局部有界”，所得的条件 (iii′)、(iv′) 或 (v′)（分别地，(iv′) 或 (v′)）仍与这些条件等价。

3.3.2. 假设 F 拟完备。设 U 是 E 的开子集，\((f_{n})\) 是从 U 到 F 的全纯映射序列，具有下列性质：

\begin{enumerate}
\def\labelenumi{(\Alph{enumi})}
\setcounter{enumi}{22}
\tightlist
\item
  U 的每一点都有一个邻域，使得序列 \((f_{n})\) 在该邻域上一致收敛。
\end{enumerate}

则序列 (f＿n) 的极限 \(f\)\((f_{n})\) 是全纯的，导数序列 \((\mathrm{D}f_{n})\)（取值于拟完备空间 \(\mathcal{L}(\mathrm{E};\mathrm{F})\)）具有性质 (W)，且 Df 是 \((\mathrm{D}f_{n})\) 的极限。

3.3.3. 设 U 是 E 的开子集，f 是从 U 到 F 的全纯映射，并假设 F 拟完备。令 \(\mathbf{R} = (\mathbf{R}_{i}) \in (\mathbf{R}_{+}^{*})^{n}\)，并假设多球 \(\mathbf{B}(\mathbf{R})\) 包含于 U 且 f 在 \(\mathbf{B}(\mathbf{R})\) 上有界。于是，对每个 \(\alpha \in N^{n}\) 以及每个 \(x = (x_{i}) \in \mathbf{B}(\mathbf{R})\)，有：

\[
\Delta^ {\alpha} f (0) (x) = \int_ {0} ^ {1} \dots \int_ {0} ^ {1} f (\mathbf {e} (\theta_ {1}) x _ {1}, \dots , \mathbf {e} (\theta_ {n}) x _ {n}) \mathbf {e} (- \alpha_ {1} \theta_ {1} - \dots - \alpha_ {n} \theta_ {n}) d \theta_ {1} \dots d \theta_ {n}
\]

\((\text{其中 } \mathbf{e}(\theta) = \exp 2\pi i \theta)\)。

此外，设 \(\gamma\) 是 \(F\) 上的连续半范数，\(M\) 是在 \(\|f(x)\|_{\gamma}\) 时 \(\|x_{i}\|=R_{i}\) 的上确界。则对所有 \(\|\Delta^{\alpha}f(0)(x)\|_{\gamma}\leqslant M\) 都有 \(x\in B(R)\)，且 \(\|\Delta^{\alpha}f(0)\|_{\gamma}\leqslant MR^{-\alpha}\)。最后，\(f\) 在 0 处的展开级数的收敛域包含多圆盘 \(B(R)\) 的内部。

3.3.4. 保持 3.3.3 的假设，并进一步假设 \(\mathbf{E}_i = \mathbf{C}\)。设 \(\sum X^{\alpha}c_{\alpha}\) 是 \(f\) 在 0 处的展开级数。有：

\[
c _ {\alpha} = \mathbf {R} ^ {- \alpha} \int_ {0} ^ {1} \dots \int_ {0} ^ {1} f (\mathbf {e} (\theta_ {1}) \mathrm{R} _ {1}, \dots , \mathbf {e} (\theta_ {n}) \mathrm{R} _ {n}) \mathbf {e} (- \alpha_ {1} \theta_ {1} - \dots - \alpha_ {n} \theta_ {n}) d \theta_ {1} \dots d \theta_ {n}
\]

以及：

\[
\| c _ {\alpha} \| _ {\gamma} \leqslant \mathbf {R} ^ {- \alpha} \sup _ {x \in \mathbf {B} (\mathbf {R})} \| f (x) \| _ {\gamma}
\]

（“柯西不等式”）。级数 \(\sum_{\alpha} X^{\alpha} c_{\alpha}\) 的严格收敛域包含 B(R) 的内部。

3.3.5. 假设 E 有限维且 F 拟完备。于是 \(\mathcal{H}(\mathbf{E};\mathbf{F})\) 中存在唯一的级数 \(f_{0}\)，它（对 E 上每个范数）具有无穷收敛半径，并满足对所有 \(f(x)=f_{0}(x)\) 都有 \(x\in E\)。

3.3.6. 若 \(f\) 是从 E 到 F 的全纯映射且 \(f(E)\) 有界，则函数 \(f\) 为常值函数（“刘维尔定理”）。

3.3.7. 假设 \(E \neq 0\)。设 f 是从 E 的开集 U 到 F 的全纯映射。设 \(a\) 是 U 中一点，\(\gamma\) 是 F 上的连续半范数。对包含于 U 的 a 的每个邻域 V，都存在 \(x \in V\)、\(x \neq a\)，使得：

\[
\| f (a) \| _ {\gamma} \leqslant \| f (x) \| _ {\gamma}.
\]

若此外 F = C，且 f 在 a 的任何邻域上都不为常值，则有 \(|f(a)| < \sup_{x \in V, x \neq a} |f(x)|\)，并且映射 f 在 a 的某邻域上是开映射。最后，设 B 是闭包包含于 U 的有界开集，B' 是其

边界。有：

\[
\sup _ {x \in \overline {{\mathbf {B}}}} | f (x) | = \sup _ {x \in \mathbf {B} ^ {\prime}} | f (x) |
\]

（“最大值原理”）。

3.3.8. 假设 E 有限维。设 U 是 E 的开子集，S 是 E 的余维 \(\geqslant2\) 的向量子空间，f 是从 \(\mathrm{U}\cap(\mathrm{E}-\mathrm{S})\) 到 F 的全纯映射。则 f 可连续延拓为从 U 到 F 的全纯映射。

假设 E = C，且 \(0 \leqslant \lambda < \mu \leqslant +\infty\)。设 f 是从开集 \(U = \{z \in C | \lambda < |z| < \mu\}\) 到 F 的全纯映射。存在唯一的 F 中元素族 \((a_{n}(f))_{n \in \mathbb{Z}}\)，使得对 U 的每个紧子集 H，族 \((a_{n}(f)z^{n})_{n \in \mathbb{Z}}\) 一致可求和，当 z 遍历 H 时其和为 \(f(z)\)（“洛朗展开”）。

再假设 \(\lambda=0\)。称 \(f\) 在点 \(x=0\) 处的阶为满足 a＿\(n\)\(a_{n}(f)\neq0\)(\(f\))0 的整数 n 的下确界（有限或无穷）。若存在 0 的邻域 V，使得 f 在 V-\(\{0\}\) 中有界，则 \(f\) 在点 \(x=0\) 处的阶为 0，并可连续延拓为开集 \(|z|<\mu\) 上的全纯函数。设 \(p\) 是整数且 p\(>0\)；若 \(f\)\(-p\) 在点 \(x=0\) 处的阶为 -\(p\)，则称 0 是 \(f\) 的 p 阶极点。

3.3.10. 假设 E = C 且 F 是赋范空间。设 f 是从 E 的开单位圆盘 U 到 F 的全纯映射，满足 \(f(0) = 0\)，并令 \(M = \sup_{z \in U} \|f(z)\|\)。于是对所有 \(\|f(z)\| \leqslant M \cdot |z|\) 都有 \(z \in U\)（“施瓦茨引理”）。

\subsection{实解析函数}

假设 K = R。

3.4.1. 设 U 是 E 的开子集，f 是从 U 到 F 的映射，并假设 F 拟完备。

下列条件等价：

\begin{enumerate}
\def\labelenumi{(\roman{enumi})}
\item
  \(f\) 是解析的。
\item
  \(f\) 属于 \(\mathbf{C}^{\infty}\) 类，并且对每个 \(a\in\mathbf{U} \)，存在 a 的邻域 \(\mathbf{V}_{a}\)\(a\) 和实数 \(M\)，使得对 F 上每个连续半范数 \(\gamma\)，存在常数 \(\mathbf{A}_{\gamma}\)，从而有
\end{enumerate}

\[
\| \Delta^ {\alpha} f (x) \| _ {\gamma} \leqslant A _ {\gamma} M ^ {| \alpha |} \quad \text {   for   every   } x \in V _ {a} \text {   and   every   } \alpha \in \mathbf {N} ^ {n}.
\]

3.4.2. 设 U 是 E 的开子集，f 是从 U 到 F 的映射。若 F 拟完备，且其强对偶 F' 是贝尔空间，则 f 解析，当且仅当对每个 \(u \circ f\)，\(u \in F'\) 都解析。

\section{超度量情形的解析函数}

本节假定 K 的绝对值是超度量的。用 \((\mathbf{E}_{i})_{1\leqslant i\leqslant n}\) 表示 K 上赋范空间的一个有限族，用 E 表示各 \(E_{i}\) 的乘积空间，并赋予范数：

\[
\| x \| = \sup \| x _ {i} \| \quad \text { if } \quad x = (x _ {i}).
\]

用 F 表示 K 上的分离多范数空间。

\subsection{收敛级数}

4.1.1. 设 \(f = \sum_{\alpha} f_{\alpha}\) 是属于 \(\hat{\mathbf{P}}(\mathbf{E}_1, \ldots, \mathbf{E}_n; \mathbf{F})\) 的形式级数（参见附录）。若 \(\gamma\) 是 \(\mathbf{F}\) 上的连续半范数，且 \(\mathbf{R} = (\mathbf{R}_i)\) 是由 \(n\) 个 \(>0\) 的实数组成的系统，置

\[
\| f \| _ {\gamma , \mathbf {R}} = \sup _ {\alpha} \| f _ {\alpha} \| _ {\gamma} \mathbf {R} ^ {\alpha}.
\]

第 3.1.1 号（第二段）和第 3.1.2 号的定义与结果不加改变地适用；特别地，定义空间

\[
\mathcal {H} _ {\mathrm{R}} (\mathrm{E} _ {1}, \dots , \mathrm{E} _ {n}; \mathrm{F}) \quad \text { and } \quad \mathcal {H} (\mathrm{E} _ {1}, \dots , \mathrm{E} _ {n}; \mathrm{F}).
\]

4.1.2. \(j\) 到 \(\hat{\mathbf{P}}(\mathbf{E};\mathbf{F})\) 的规范同构 \(\hat{\mathbf{P}}(\mathbf{E}_{1},\ldots,\mathbf{E}_{n};\mathbf{F})\) 通过限制，对每个 \(\mathcal{H}_{\mathbb{R}}(\mathbf{E};\mathbf{F})\)，给出从 \(\mathcal{H}_{(\mathbb{R},\ldots,\mathbb{R})}(\mathbf{E}_{1},\ldots,\mathbf{E}_{n};\mathbf{F})\) 到 \(\mathbf{R}\in\mathbb{R}_{+}^{*}\) 的拓扑向量空间同构；它还给出从 \(\mathcal{H}(\mathbf{E};\mathbf{F})\) 到 \(\mathcal{H}(\mathbf{E}_{1},\ldots,\mathbf{E}_{n};\mathbf{F})\) 的同构。更确切地说，若 \(f=\sum_{m}f_{m}\in\hat{\mathbf{P}}(\mathbf{E};\mathbf{F})\) 且 \(j(f)=\sum_{\alpha}f_{\alpha}\)，则对 \(\gamma\) 上的每个连续半范数 \(\mathbf{F}\)，有：

\[
\begin{array}{l} \| f _ {m} \| _ {\gamma} = \sup _ {| \alpha | = m} \| f _ {\alpha} \| _ {\gamma} \\ \| f \| _ {\gamma , \mathbf {R}} = \| j (f) \| _ {\gamma , (\mathbf {R}, \dots , \mathbf {R})}. \end{array}
\]

4.1.3. 设 \(f = \sum_{\alpha} f_{\alpha}\) 是 \(\mathcal{H}(E_1, \ldots, E_n; F)\) 中的元素；令 \(I(f)\) 为满足如下条件的 \(R \in (\mathbf{R}_+^*)^n\) 的集合：对 F 上每个连续半范数 \(\gamma\)，当 \(\|f_{\alpha}\|_{\gamma} R^{\alpha}\) 趋于无穷时，乘积 \(|\alpha|\) 趋于零。集合 \(I(f)\) 非空；称其为 \(f\) 的严格收敛指标。点集 \(\Omega(f)\)

\[
(\log \mathbf {R} _ {1}, \dots , \log \mathbf {R} _ {n}) \quad \text { for } \quad \mathbf {R} \in I (f)
\]

是 \(\mathbf{R}^{n}\) 的凸子集。

当 \(n=1\) 时，集合 \(\mathbf{I}(f)\) 是 \(\mathbf{R}\) 中左开且右端开或闭的区间；其上界（有限或 \(+\infty\)）记作 \(\rho(f)\)，称为 \(f\) 的严格收敛半径。

使用 4.1.2 的记号，有 \(\mathbf{R} \in \mathbf{I}(f)\) 当且仅当

\[
(\mathbf {R}, \dots , \mathbf {R}) \in \mathrm{I} (j (f)).
\]

满足存在 \(x = (x_i)\) 且对 \(\mathbf{R} = (\mathbf{R}_i) \in I(f)\) 有 \(\| x_i \| \leqslant \mathbf{R}_i\) 的点 \(1 \leqslant i \leqslant n\) 的集合称为 \(f\) 的严格收敛域，记作 \(C(f)\)。它是 \(E\) 的开子集，是多球

\[
\mathrm{B} (\mathbf {R}) = \{x \in \mathrm{E} | \| x _ {i} \| \leqslant \mathrm{R} _ {i} \text {   for   } 1 \leqslant i \leqslant n \},
\]

当 \(\mathbf{R}\in\mathbf{I}(f)\) 时的并。

4.1.4. 将其中所有的 \(\tilde{\mathbf{C}}(f)\) 换为 \(\mathbf{C}(f)\)，将 \(\tilde{\mathcal{H}}_{\mathbb{R}}\) 换为 \(\mathcal{H}_{\mathbb{R}}\) 后，3.1.7 和 3.1.8 的结果仍然成立。

4.1.5. 设 \(\mathbf{F}_1, \ldots, \mathbf{F}_m\) 是完备赋范空间，且 \(\mathbf{F}\) 拟完备。令 \(\mathbf{f} = (f_i)_{1 \leqslant i \leqslant m}\)，其中 \(f_i \in \mathcal{H}(\mathbf{E}_1, \ldots, \mathbf{E}_n; \mathbf{F}_i)\)，并令 \(g \in \mathcal{H}(\mathbf{F}_1, \ldots, \mathbf{F}_m; \mathbf{F})\)，假设 \((f_i(0))_{1 \leqslant i \leqslant m}\) 中的点 \(\mathbf{E}\) 属于 \(g\) 的严格收敛域。于是，对每个 \(\alpha \in \mathbb{N}^m\)，形式级数 \(g_\alpha \circ \mathbf{f}\) 属于 \(\mathcal{H}(\mathbf{E}_1, \ldots, \mathbf{E}_n; \mathbf{F})\)，且族 \(g_\alpha \circ \mathbf{f}\) 在 \(\mathcal{H}(\mathbf{E}_1, \ldots, \mathbf{E}_n; \mathbf{F})\) 中可求和（从而当然也在 \(\hat{\mathbf{P}}(\mathbf{E}_1, \ldots, \mathbf{E}_n; \mathbf{F})\) 中可求和）。其和记作 \(g \circ \mathbf{f}\)。

更确切地说，存在 \(\mathbb{R} \in \bigcap_{i} \mathrm{I}(f_i)\) 和 \(\mathbb{R}' \in \mathrm{I}(g)\)，使得对 \(\sup_{|\alpha| > 0} \| f_{i,\alpha} \| \mathbb{R}^{\alpha} < \mathbb{R}_i'\) 有 \(1 \leqslant i \leqslant m\)。在这些条件下，形式级数 \(g_{\alpha} \circ \mathbf{f}\) 属于 \(\mathcal{H}_{\mathbb{R}}(\mathbf{E}_1, \ldots, \mathbf{E}_n; \mathbf{F})\)，且族 \(g_{\alpha} \circ \mathbf{f}\) 在 \(\mathcal{H}_{\mathbb{R}}(\mathbf{E}_1, \ldots, \mathbf{E}_n; \mathbf{F})\) 中可求和。最后，若 \(x \in \mathbf{B}(\mathbb{R})\)，则 \(\mathbf{f}(x) = (f_i(x))\) 属于 \(\mathbf{C}(g)\)，且有：

\[
g (\mathbf {f} (x)) = (g \circ \mathbf {f}) (x).
\]

再假设，对每个 \(i\)，存在 E＿i 中元素组成的族 \((e_{j}^{i})\)，使得 \(\mathbf{E}_{i}\) 中每个元素 \(x\)\(\mathbf{E}_{i}\) 都是可求和族 \((\lambda_{j}e_{j}^{i})\) 的和（其中 \(\lambda_{j}\in\mathbf{K}\)），且 \(\|x\|=\sup_{j}|\lambda_{j}|\)。于是，在前一段中可将条件 \(\sup_{|\alpha|>0}\|f_{i,\alpha}\|\mathbf{R}^{\alpha}<\mathbf{R}_{i}^{\prime}\) 换为条件 \(\|f_{i}\|_{\mathbf{R}}\leqslant\mathbf{R}_{i}^{\prime}\)。

4.1.6. 假设 \(E_{i} = K\)（\(1 \leqslant i \leqslant n\)）。于是将空间 \(\hat{\mathbf{P}}(\mathbf{K}^{n}; \mathbf{F})\) 识别为具有 F 中系数的 n 个不定元 \(X_{1}, \ldots, X_{n}\) 的形式幂级数空间，且 \(\hat{\mathbf{P}}(\mathbf{K}^{n}; \mathbf{F})\) 中的元素写作：

\[
f = \sum_ {\alpha} X ^ {\alpha} c _ {\alpha} \quad \text { with } c _ {\alpha} \in F.
\]

\[
\| f \| _ {\gamma , \mathbf {R}} = \sup \| c _ {\alpha} \| _ {\gamma}. \mathbf {R} ^ {\alpha}.
\]

若 \(\mathbf{R} \in (\mathbf{R}_{+}^{*})^{n}\)，且 \(\gamma\) 是 \(F\) 上的连续半范数，则有：

3.1.10 的最后一段仍然成立，3.1.11 也同样成立。

4.1.7. 假设 K 是完备赋值域（必为超度量的）L 的闭子域。对于 \(y \in E_{i} \otimes_{K} L\)，置：

\[
\| y \| = \inf (\sup _ {k} | a _ {k} |. \| x _ {k} \|)
\]

下确界取遍满足 \((x_{k}, a_{k}) \in E_{i} \times L\) 的所有有限对族 \(y = \sum_{k} x_{k} \otimes a_{k}\)。由此在 L-向量空间 \(E_{i} \otimes_{K} L\) 上得到一个范数，它在 \(E_{i}\) 上诱导出给定范数。用 \(E_{i}^{L}\) 表示 \(E_{i} \otimes_{K} L\) 关于此范数的完备化。

现在设 F 是分离完备的 L-向量多范数空间。对于 \(f_{\alpha}\) 上取值于作为 F 底层空间的 K-向量多范数空间 \(\alpha\) 的每个 K-多项式连续映射 \(E_1 \times \cdots \times E_n\)，若其关于多重次数 \(F_K\) 齐次，则存在唯一的在 \(\tilde{f}_{\alpha}\) 上取值于 F 的 L-多项式连续映射 \(E_1^L \times \cdots \times E_n^L\)，关于同一多重次数齐次，并延拓 \(f_{\alpha}\)。对 L-向量空间 F 上的每个连续半范数，有

\[
\| \tilde {f} _ {\alpha} \| _ {\gamma} = \| f _ {\alpha} \| _ {\gamma}.
\]

若 \(f = \sum_{\alpha} f_{\alpha} \in \mathcal{H}(E_1, \ldots, E_n; F_K)\)，则 \(\tilde{f} = \sum_{\alpha} \tilde{f}_{\alpha} \in \mathcal{H}(E_1^L, \ldots, E_n^L; F)\)。级数 \(f\) 和 \(\tilde{f}\) 具有相同的严格收敛指标（且当 \(n = 1\) 时具有相同的严格收敛半径）。

反之，令 L 为 K 的非离散闭子域，令 \(E_{i}^{0}\) 和 \(F^{0}\) 为分别从 \(E_{i}\) 和 F 限制标量得到的 L 上的空间。若 \(f = \sum f_{\alpha} \in \mathcal{H}(E_{1}, \ldots, E_{n}; F)\)，则 \(f_{\alpha} \in \mathrm{P}_{\alpha}(E_{1}^{0}, \ldots, E_{n}^{0}; F^{0})\)；若置 \(f^{0} = \sum_{\alpha} f_{\alpha} \in \hat{\mathrm{P}}(E_{1}^{0}, \ldots, E_{n}^{0}; F^{0})\)，则 \(f^{0} \in \mathcal{H}(E_{1}^{0}, \ldots, E_{n}^{0}; F^{0})\)。有 \(C(f) \subset C(f^{0})\)，且对所有 \(f(x) = f^{0}(x)\)，有 \(x \in C(f)\)。

\subsection{解析函数}

4.2.1. 3.2.1 和 3.2.2 的定义与结果不加改变地成立。

4.2.2. 使用 3.2.2 的记号，\(\Delta^{a}f\) 在 U 中一点 \(a\) 处的幂级数展开的严格收敛指标包含 \(f\) 在 \(a\) 处的幂级数展开的严格收敛指标。

4.2.3. 3.2.4、3.2.5、3.2.7、3.2.8 和 3.2.11 的结果仍准确成立。若再假设 K 的特征为零，3.2.6 的结果也成立。



4.2.4. 假设 F 拟完备，且 \(f\in\mathcal{H}(E_{1},\ldots,E_{n};F)\)。函数 \(x\mapsto f(x)\) 在 C(f) 中解析。对每个 \(a\in\mathbf{C}(f)\)，\(f\) 在 \(a\) 处的幂级数展开的收敛指标等于 \(f\) 的收敛指标。

\subsection{若干不等式}

4.3.1. 假设 K 至少满足下列条件之一：(a) K 的剩余域是无限的；

\begin{enumerate}
\def\labelenumi{(\alph{enumi})}
\setcounter{enumi}{1}
\tightlist
\item
  K 在映射 \(a \mapsto |a|\) 下的像在 \(\mathbf{R}_{+}\) 中稠密。（换言之，假设 K 不是局部紧的。）
\end{enumerate}

设 \(f = \sum_{\alpha} f_{\alpha} \in \mathcal{H}(E_1, \ldots, E_n; F)\)，且 \(R \in I(f)\)。有：

\[
\sup _ {x \in \mathbf {B} (\mathbb {R})} \| f (x) \| _ {\gamma} = \sup _ {\alpha} \sup _ {x \in \mathbf {B} (\mathbb {R})} \| f _ {\alpha} (x) \| _ {\gamma}
\]

对 F 上每个连续半范数 γ 均成立（“柯西不等式”）。

4.3.2. 存在常数 \(a > 0\)，使得对每个连续齐次多项式 \(f_{\alpha} \in \mathbf{P}_{\alpha}(\mathrm{E}_{1}, \ldots, \mathrm{E}_{n}; \mathbf{F})\) 以及每个 \(\mathbf{R} \in (\mathbf{R}_{+}^{*})^{n}\)，有：

\[
a ^ {| \alpha |} \mathbf {R} ^ {\alpha} | \alpha ! | \| f _ {\alpha} \| _ {\gamma} \leqslant \sup _ {x \in \mathbf {B} (\mathbf {R})} \| f _ {\alpha} (x) \| _ {\gamma} \leqslant \| f _ {\alpha} \| _ {\gamma} \mathbf {R} ^ {\alpha}
\]

对 F 上每个连续半范数 \(\gamma\) 均成立。若 K 满足 4.3.1 的条件 (b)，或者若 \(E_{i}\) 在映射 \(x \mapsto \|x\|\) 下的像包含于 K 在映射 \(a \mapsto |a|\) 下的像中且包含 \(\mathbf{R}_{i}\)（\(1 \leqslant i \leqslant n\)），则可取 \(a = 1\)。

4.3.3. 若 K 的特征为零，则形式级数 \(f = \sum_{\alpha} f_{\alpha}\) 属于 \(\mathcal{H}(E_1, \ldots, E_n; F)\)，当且仅当存在 \(R \in (R_+^*)^n\)，使得

\[
\sup _ {\alpha} \sup _ {x \in B (R)} \| f _ {\alpha} (x) \| _ {\gamma} <   + \infty
\]

对 F 上每个连续半范数 γ 均成立。

\section{流形}

\subsection{坐标图、图册与流形}

5.1.1. 设 X 是一个集合。X 上的一个坐标图是三元组 \(c = (U, \varphi, E)\)，其中 U 是 X 的子集，E 是巴拿赫空间，\(\varphi\) 是从 U 到 E 的开子集的双射。称 U 为坐标图 c 的定义域。若 E 的维数为 n，则称 c 的维数为 n。否则置 \(\dim c = +\infty\)。

5.1.2. 当不会对 r 产生歧义时，若 X 的两个坐标图 \(c = (\mathbf{U}, \varphi, \mathbf{E})\) 和 \(c' = (\mathbf{U}', \varphi', \mathbf{E}')\)\(X\) 满足下列条件，则称它们 \(C^r\)\(r\)-相容（或简写为相容）：

\begin{enumerate}
\def\labelenumi{(\alph{enumi})}
\item
  \(\varphi (\mathbf{U}\cap \mathbf{U}^{\prime})\)（分别地，\(\varphi^{\prime}(\mathbf{U}\cap \mathbf{U}^{\prime}))\)）在 \(\mathbf{E}\)（分别地，\(\mathbf{E}'\)）中是开集；
\item
  从 \(\varphi \circ \varphi'^{-1}\) 到 \(\varphi' \circ \varphi^{-1}\) 的映射 \(\varphi'(U \cap U')\)（分别地，从 \(\varphi(U \cap U')\) 到 \(\varphi(U \cap U')\) 的映射 \(\varphi'(U \cap U')\)）属于 \(C^r\) 类（参见 2.3.1、3.2.1 和 4.2.1）。
\end{enumerate}

5.1.3. 称集合 X 的一个 \(\mathbf{C}^{r}\)-图册（或简写为图册）为 X 的一族坐标图，若这些坐标图两两 \(\mathbf{C}^{r}\)-相容且其定义域的并为 X。若 \(\mathbf{C}^{r}\) 是图册，则称 X 的两个图册 A 和 B \(A \cup B\)-等价。图册之间的 \(\mathbf{C}^{r}\)-等价关系是等价关系。

5.1.4. 设 \(\mathfrak{S}\) 是巴拿赫空间的集合。若图册 \(\mathcal{A}\) 的每个坐标图 c = (U, , E) 都满足 E  \(\mathfrak{S}\)，则称图册 \(E \in \mathfrak{S}\) 为 \(c = (U, \varphi, E)\) 型。同样，若图册 \(\mathcal{A}\) 的每个坐标图 \(\mathcal{A}\) 中的 \(E\)\((U, \varphi, E)\) 都是希尔伯特空间（分别地，具有可数型的希尔伯特空间），则称 \(\mathcal{A}\) 为希尔伯特型（分别地，具有可数维的希尔伯特型）。

5.1.5. 称一个 \(\mathbf{C}^{r}\) 类 K-流形（或 K 上的 \(\mathbf{C}^{r}\) 类流形；当 K 和 r 不会引起歧义时也简写为流形）为赋予了关于 \(\mathbf{C}^{r}\)-等价关系的一类等价图册的集合 X（Sets，第 II 章，第 6 节，第 9 号）。此类中的图册称为流形 X 的图册。属于 X 的一个图册的坐标图称为流形 X 的坐标图。定义域包含点 a \(\in X\) 的 X 的坐标图称为 X 在 a 处的坐标图。以 a 为中心的坐标图是 X 在 a 处的坐标图 \(\varphi\)，满足 \(\varphi(a) = 0\)。

若 X 是集合，A 是 X 的图册，则赋予 A 的等价类的集合 X 称为由 A 定义的流形。

当 \(r \neq \omega\)（因而 \(K = R\)）时，\(C^r\) 类流形有时称为可微流形。\(C^\omega\) 类流形也称为 \(K\) 上的解析流形（或 \(K\)-解析流形）。

此外，当 K = R（分别地，C、\(Q_{p}\)）时，也称实解析流形（分别地，复解析流形、p-进解析流形）。

5.1.6. 设 X 是流形。X 中作为 X 的坐标图定义域之并的子集构成 X 上某个拓扑的开集族，该拓扑称为流形结构的底层拓扑。对于 X 的每个坐标图 \(c = (U, \varphi, E)\)，映射 \(\varphi\) 是从赋予 X 所诱导拓扑的开集 U 到 E 的开集 \(\varphi(U)\) 的同胚。

X 的底层拓扑空间是贝尔空间。当 K 等于 R 或 C 时，它是局部连通的。

设 X 是流形，A 是 X 的图册。拓扑空间 X 为豪斯多夫空间的充要条件是满足下列条件：对于图册 A 中任意坐标图 (U，\(\varphi\)，E) 和 (V，\(\psi\)，F)，从 \(\psi \circ \varphi^{-1}\) 到 \(\varphi(\mathrm{U} \cap \mathrm{V})\) 的映射 \(\psi(\mathrm{U} \cap \mathrm{V})\) 的图像在 \(\varphi(\mathrm{U}) \times \psi(\mathrm{V})\) 中是闭的。

设 X 是流形；假设拓扑空间 X 是正则的。则对每个点 \(a \in X\)，存在 X 在 a 处的一个坐标图 (U，\(\varphi\)，E)，具有如下性质：U 的子集 Y 在 X 中闭的充要条件是其像 \(\varphi(Y)\) 在 E 中闭。若空间 X 是仿紧的，则 X 上存在一个与 X 的拓扑相容的距离，并使 X 成为完备度量空间。

5.1.7. 设 S 是巴拿赫空间的集合。若一个流形具有 S 型（分别地，希尔伯特型、可数维希尔伯特型）图册，则称其为 S 型流形（分别地，希尔伯特型流形、可数维希尔伯特型流形）。若 S 约化为单个元素 E，则 S 型流形也称为 E 型纯流形。

维数为 n 的纯流形称为 \(K^{n}\) 型纯流形。若流形是 S 型的，其中 \(S = \{K^{n}; n \in N\}\)，则称其局部有限维。

5.1.8. 设 X 是流形，且 \(x \in X\)。X 在 x 处的坐标图的维数（有限或 \(+\infty\)）只取决于 x。称其为 X 在 x 处的维数，记作 \(\dim_{x}X\)。X 的维数记作 \(\dim X\)，是 \(\dim_{x}X\) 时 \(x \in X\) 的上界。

函数 \(x \mapsto \dim_x X\) 局部为常值。流形 \(X\) 局部有限维（分别地，纯 \(n\) 维）的充要条件是对所有 \(\dim_x X < +\infty\) 都有 \(\dim_x X = n\)（分别地，\(x \in X\)）。

5.1.9. 假设 K 局部紧。设 X 是分离流形。X 局部有限维的充要条件是 X 局部紧。

5.1.10. 设 X 是流形，\(\xi^{1},\ldots,\xi^{n}\) 是从 X 的子集 U 到 K 的映射。若三元组 (U，\(\xi=(\xi^{1},\ldots,\xi^{n})\)，K \(\xi\)) 是 X 的坐标图，则称 \(^{n}\) 是 X 在 U 中的坐标系；该坐标图也记作 (U；\(\xi\)) 或 (U；\(\xi^{1},\ldots,\xi^{n}\))。若 \(a\in\mathbf{U}\)，也称 \(\xi\) 是 X 在 a 处的坐标系；若进一步对 \(\xi^{i}(a)=0\) 有 \(i=1,2,\ldots,n\)，则称坐标系 \(\xi\) 以 a 为中心。

\subsection{流形的例子}

5.2.1. 设 X 是集合；X 上存在唯一的流形结构，使其底层拓扑空间是离散的；该结构是维数为 0 的纯流形结构。

5.2.2. 设 E 是巴拿赫空间。三元组 \(c = (\mathrm{E}, \mathrm{Id}_{\mathrm{E}}, \mathrm{E})\) 是 E 的坐标图，且 \(\mathcal{A} = \{c\}\) 是 E 的图册，因而在 E 上定义 E 型纯流形结构；底层拓扑是 E 上给定的拓扑。以下谈及 E 上的流形结构时，始终指上述结构。

特别地，这适用于每个赋予与其向量结构相容的唯一豪斯多夫拓扑的 K 上有限维向量空间（Esp. Vect. Top.，第 I 章，第 2 节，第 3 号）。

5.2.3. 设 X 是流形，U 是 X 的开子集。U 上存在一种流形结构，其坐标图是定义域包含于 U 的流形 X 的坐标图。称此结构由 X 的结构诱导（参见第 5.3.4 号）；赋予此结构的 U 称为 X 的开子流形。

特别地，巴拿赫空间 E 的每个开子集都赋有 E 型纯流形的规范结构。设 X 是流形；三元组 (U，\(\varphi\)，E) 是 X 的坐标图的充要条件是 U 在 X 中开，且 \(\varphi\) 是从 X 的开子流形 U 到 E 的开子流形的同构。

5.2.4. 设 X 是集合，\((X_{i})_{i\in I}\) 是 X 的一个覆盖。假设每个 \(X_{i}\) 上给定一种流形结构，满足下列条件：

对 I 中任意 \(i\) 和 \(j\)，集合 \(\mathbf{X}_{i} \cap \mathbf{X}_{j}\) 在 \(\mathbf{X}_{i}\) 和 \(\mathbf{X}_{j}\) 中都是开集，并且由 \(\mathbf{X}_{i} \cap \mathbf{X}_{j}\) 与 X＿j 的结构在 \(\mathbf{X}_{i}\)  \(\mathbf{X}_{j}\) 上诱导的流形结构相同。

于是 X 上存在唯一的流形结构，使得对 I 中每个 i，\(X_{i}\) 都是 X 的开子流形。称流形 X 由各流形 \(X_{i}\) 粘合而成。

5.2.5. 设 X 是流形。X 中满足 \(X_{n}\) 的点 x 的集合 \(\dim_{x}X=n\ (n\ \text{为整数}\geqslant0)\) 是 X 的开子流形，且是纯 n 维的。

5.2.6. 设 E 是巴拿赫空间。可以按如下方式在具有拓扑补的 E 的向量子空间集合 \(\mathbf{G}(\mathbf{E})\) 上赋予解析流形结构：对每一对 \((\mathrm{F}_0, \mathrm{G}_0) \in \mathbf{G}(\mathbf{E}) \times \mathbf{G}(\mathbf{E})\)，若 \(\mathrm{E} = \mathrm{F}_0 \oplus \mathrm{G}_0\)，用 \(\mathrm{U}_{\mathrm{G}_0}\) 表示具有 \(\mathrm{F} \in \mathbf{G}(\mathbf{E})\) 作为补空间的 \(\mathrm{G}_0\) 的集合，并通过将每个 \(\varphi_{\mathrm{F}_0, \mathrm{G}_0}\) 与从 \(\mathrm{U}_{\mathrm{G}_0}\) 到 \(\mathcal{L}(\mathrm{F}_0; \mathrm{G}_0)\) 的映射关联起来定义从 \(\mathrm{F} \in \mathrm{U}_{\mathrm{G}_0}\) 到巴拿赫空间 \(\mathrm{F}_0\) 的双射 \(\mathrm{G}_0\)，该映射的图像是 \(\mathrm{E} = \mathrm{F}_0 \times \mathrm{G}_0\) 的子空间 F。坐标图 \((\mathrm{U}_{\mathrm{G}_0}, \varphi_{\mathrm{F}_0, \mathrm{G}_0}, \mathcal{L}(\mathrm{F}_0; \mathrm{G}_0))\) 构成 \(\mathbf{G}(\mathbf{E})\) 的一个图册。赋予由此图册定义的流形结构后，\(\mathbf{G}(\mathbf{E})\) 称为 E 的格拉斯曼流形。

拓扑空间 G(E) 可度量化。若 K 局部紧且 E 有限维，则 G(E) 紧。

对每个整数 \(n \geqslant 0\)，由维数（分别地，余维）为 \(\mathbf{G}_{n}(\mathbf{E})\) 的 \(\mathbf{G}^{n}(\mathbf{E})\) 构成的集合 \(\mathbf{F} \in \mathbf{G}(\mathbf{E})\)（分别地，G＾\(n\)(E)）在 \(\mathbf{G}(\mathbf{E})\) 中既开又闭，并且是 \(\mathbf{G}(\mathbf{E})\) 的纯开子流形。若 \(\mathbf{K} = \mathbf{R}\) 或 \(\mathbf{C}\)，则对每个 n，\(\mathbf{G}_{n}(\mathbf{E})\)\(n\) 都是连通的。

将每个 \(F \in \mathbf{G}(\mathbf{E})\) 关联到其在 E 的对偶 \(E'\) 中的正交补的映射，是从 \(\mathbf{G}(\mathbf{E})\) 到 \(\mathbf{G}(\mathbf{E}')\) 的态射，并对每个整数 n 诱导出从 \(\mathbf{G}^{n}(\mathbf{E})\) 到 \(\mathbf{G}_{n}(\mathbf{E}')\) 的同构。

若 K = R 或 C，或者 E 有限维，则 \(\mathbf{G}_{1}(\mathbf{E})\) 正是与 E 关联的射影空间，因而赋有解析流形结构。

\subsection{\texorpdfstring{\(C^{r}\) 类函数与流形态射}{C\^{}\{r\} 类函数与流形态射}}

5.3.1. 设 X 是 \(\mathbf{C}^{r}\) 类流形，F 是分离多范数空间，且 \(f\) 是从 X 到 F 的映射。若对于 X 的每个坐标图 \(f\)，从 \(\mathbf{C}^{r}\) 到 \((V, \varphi, E)\) 的映射 \(f\circ\varphi^{-1}\) 属于 \(\varphi(\mathrm V)\) 类，则称 \(F\) 属于 \(\mathbf{C}^{r}\) 类。只需对 X 的一个图册中的坐标图满足此条件即可。从 X 到 F 的 \(\mathbf{C}^{r}\) 类映射构成所有从 X 到 F 的映射空间的向量子空间，记作 \(\mathcal{C}^{r}(\mathrm{X};\mathrm{F})\)。

当 F = K 时，置 \(\mathcal{C}^{r}(X;K)=\mathcal{C}^{r}(X)\)；它是从 X 到 K 的映射代数的 K-子代数。\(\mathcal{C}^{r}(X)\) 的元素也称为态射函数。

当 X 是巴拿赫空间的开子集时，此术语与第 2.3.1、3.2.1 和 4.2.1 号中的术语一致。

5.3.2. 设 X 和 Y 是两个 \(\mathbf{C}^{r}\) 类流形，\(f\) 是从 X 到 Y 的映射。若 \(f\)\(\mathbf{C}^{r}\)\(\mathbf{C}^{r}\) 连续，且对 Y 的每个坐标图 \((\mathbf{V},\psi ,\mathbf{F})\)

则称 f 属于 \(\psi \circ f\)\(f^{-1}(\mathbf{V})\) 类，或称为（\(\mathbf{C}^{r}\) 类）流形态射。只需存在 Y 的一个图册 \(\mathcal{A}\)，使得对每个坐标图 \((\mathbf{V}, \psi, \mathbf{F}) \in \mathcal{A}\)，集合 \(f^{-1}(\mathbf{V})\) 在 X 中开，且从 X 的开子流形 \(\psi \circ f\) 到 F 的映射 \(f^{-1}(\mathbf{V})\) 属于 \(\mathbf{C}^{r}\) 类。X 到 Y 的态射集合记作 \(\mathcal{C}^{r}(\mathbf{X}; \mathbf{Y})\)。当 Y 是赋予其规范流形结构的巴拿赫空间时，5.3.1 和 5.3.2 的定义相容。属于 \(\mathbf{C}^{\omega}\) 类的映射也称为 K-解析映射（或简写为解析映射）。当 \(\mathbf{K} = \mathbf{C}\) 时，也称为全纯映射。

设 (U，\(\varphi\)，E) 是 X 的坐标图，(V，\(\psi\)，F) 是 Y 的坐标图，且 \(f(\mathbf{U}) \subset \mathbf{V}\)。从 E 的开子集 \(\psi \circ f \circ \varphi^{-1}\) 到 F 的开子集 \(\varphi(\mathbf{U})\) 的映射 \(\psi(\mathbf{V})\) 称为 \(f\) 在给定坐标图中的表达式。

5.3.3. 假设 X 和 Y 有限维；设 \(a \in \mathbf{X}\)、\(b \in \mathbf{Y}\)，且 \(f\) 是从 X 到 Y 的映射，满足 \(f(a) = b\)。首先假设 \(f\) 属于 \(\mathbf{C}^r\)\((\mathrm{U}; \xi^1, \ldots, \xi^m)\) 类，并分别考虑 X 在 \(a\) 处的坐标系 \((\mathrm{V}; \eta^1, \ldots, \eta^n)\) 和 Y 在 \(b\) 处的坐标系 (V; ＾1, , ＾n)，满足 \(f(\mathrm{U}) \subset \mathrm{V}\)。令 \(\xi\) 为从 U 到 \((\xi^1, \ldots, \xi^m)\) 的映射 \(\mathbf{K}^m\)。于是存在 \(u^j\) 的开子集 \(\mathbf{C}^r\) 上取值于 K 的 \(\xi(\mathrm{U})\) 类函数 \(\mathbf{K}^m\)，使得 Y 中点 \(y = f(x)\) 的坐标（\(x\) 在 U 中）由下列公式给出：

\[
\eta^ {j} (y) = u ^ {j} (\xi^ {1} (x), \dots , \xi^ {m} (x)) \quad \text { for } 1 \leqslant j \leqslant n\tag{1}
\]

这等价于：

\[
\eta^ {j} \circ f = u ^ {j} (\xi^ {1}, \dots , \xi^ {m}) \quad \text { for } 1 \leqslant j \leqslant n.\tag{2}
\]

称上述公式为借助所选坐标系得到的 f 的表达式。

反之，若对 X 的每个点 \(a\)，都能找到 X 在 \(a\) 处的坐标系和 Y 在 \(b = f(a)\) 处的坐标系，使其满足上述条件，则 \(f\) 属于 \(\mathbf{C}^{r}\) 类。

5.3.4. 流形态射满足 Ens. 第 IV 章第 2 节的公理：

\begin{enumerate}
\def\labelenumi{\alph{enumi})}
\item
  两个态射的复合是态射。
\item
  双射 \(f:X\to Y\) 是同构的充要条件是 \(f\) 和 \(f^{-1}\) 都是态射。
\end{enumerate}

5.3.5. 假设 \(r = \omega\)\(f, g\)。设 X 是流形，\(f\)、\(g\) 是从 X 到分离多范数空间或分离解析流形的两个解析映射。存在某邻域使 f 和 g 相等的点的集合，在 X 中既开又闭。

5.3.6. 假设 K = R 且 \(r \neq \omega\)。若对 X 的每个局部有限开覆盖，都存在从属于该覆盖且由 \(C^{r}\) 类函数组成的连续单位分解（Top. Gén.，第 IX 章，第 4 节，第 3 号），则称流形 X 具有 \(C^{r}\) 类单位分解。

设 E 是巴拿赫空间；考虑下列性质：

对于 E 的任意两个不相交闭子集 A 和 B，存在 E 上取值于 R 的 \(C^{r}\) 类函数 f，使得对 \(f(x) = 0\) 有 \(x \in A\)，对 \(f(x) = 1\) 有 \(x \in B\)，且对所有 \(0 \leqslant f(x) \leqslant 1\) 有 \(x \in E\)。

若 S 是具有性质 (PU) 的巴拿赫空间集合，则每个 S 型仿紧流形都具有 \(\mathbf{C}^{r}\) 类单位分解。

每个有限维空间和每个可分希尔伯特空间都具有性质 (PU)。

5.3.7. 假设 K 是超度量的。设 X 是仿紧流形。对 X 的每个开覆盖 U，都存在从属于覆盖 U 的 X 的开闭子集划分。

5.3.8. 假设 K 局部紧。设 X 和 Y 是两个有限维纯流形，f 是从 X 到 Y 的态射。若 dim X \textless{} dim Y，且 X 的拓扑具有可数开集基，则 \(f(X)\) 在 Y 中是第一纲的。
